% SPDX-FileCopyrightText: 2026 Will Cook
% SPDX-License-Identifier: CC-BY-4.0
%
% Long-form reasoning record seeded from the complete problem note.
% The short note and this record are now independent publication units.
% This flat manuscript is generated from reviewable parts below.
\documentclass[11pt]{article}
\input{problem-note-preamble}
\renewcommand{\commit}{85054358cdca543531b473684e9c9d1c0ece6dd7}
\renewcommand{\commitshort}{85054358cdca}
\usepackage{longtable}
\hypersetup{
  pdftitle={The Factorial-Denominator Series: complete reasoning record for Erdős Problem 68},
  pdfsubject={Erdős Problem Notes: Problem 68},
  pdfkeywords={irrationality, factorial denominators, channel radius, moving private factors, companion orbit, strict successor, Lean 4},
}

% The prose keeps compact declaration labels; the source appendix supplies the
% commit-pinned links checked by the public release gate.
\newcommand{\pdecl}[1]{\emph{\leanlabel{#1}}}
\newcommand{\Ser}{S}
\newcommand{\chn}[2]{V_{#1}(#2)}
\newcommand{\Lc}{L}
\newcommand{\den}{\operatorname{den}}

\runninghead{Erd\H{o}s \#68: complete reasoning record}
\title{The Factorial-Denominator Series: Complete Reasoning Record}
\subtitle{Erd\H{o}s Problem~\#68}
\author{Will Cook}
\date{1 September 2026}

\setlength{\emergencystretch}{2em}
\def\UrlBreaks{\do\.\do\_\do\/}
\begin{document}
% ---- part core ----
\maketitle
\noteprovenance

\begin{abstract}
\noindent
Write $d_n=n!-1$, $H_M=\sum_{2\le n\le M}d_n^{-1}$ and
$\Lc_M=\lcm(d_2,\ldots,d_M)$.  For a prime $p$ with $e=v_p(\Lc_M)\ge1$, let
$J=\{n:2\le n\le M,\ v_p(d_n)=e\}$ and write $d_n=p^eu_n$ on $J$.  Then
\[
 v_p\bigl(\den(H_M)\bigr)=e
 \quad\Longleftrightarrow\quad
 \sum_{n\in J}u_n^{-1}\ne0\ \hbox{ in }\mathbb F_p,
\]
the inverses being taken modulo $p$.  The prefixes through $138$ and $2592$
cancel the primes $139$ and $2593$ completely.  Second, the common
denominator itself grows fast:
\[
 \liminf_{N\to\infty}\frac{\log \Lc_N}{N^{3/2}\log N}\ \ge\ \frac{2\sqrt2}{3},
\]
by an elementary segment argument that uses
$\gcd(i!-1,j!-1)\mid j!/i!-1$ and no external multiplicity theorem.  The two
statements sit either side of reduction, and the cancellations at $139$ and
$2593$ are why the second cannot be substituted for the first.

Wilson's theorem gives arbitrarily late first occurrences of primes among the
factorial gaps, so prefix-private prime support of the denominators occurs
cofinally, with no bound relating the prime to its first hit.  Four exact
criteria locate the irrationality of $\Ser=\sum_{n\ge2}d_n^{-1}$: a
strict-successor carry equivalence, a companion-orbit residue condition, a
lower-interval escape condition of width $O(m^{-2})$, and a boundary covering
the whole shifted family $\sum_{n\ge2}1/(n!+t)$ for every integer $t\ge-1$,
whose member $t=0$ returns the irrationality of $e$.  A finite-channel radius
bound gives $3t^3<2(R+1)$ under the exact cancellation and factorial-size
hypotheses, and the growth statement above raises the asymptotic constant in
that estimate from $3/2$ to $16/9$.  Two finite computations exclude
denominators: $q\nmid299999!$ from an exact interval carry census through
$300000$, and the kernel-checked bounds $q\ge2^{39990}$ and
$q>10^{12040}$ from a certified continued-fraction enclosure. The
factorial-grid exclusion uses the external GMP computation and remains a
separate outstanding Lean obligation.  The irrationality of $\Ser$ is open, and the remaining
arithmetic inputs are stated in the last section.
\end{abstract}

\section{Which prime powers survive reduction}\label{long68:sec:prime-powers}

\begin{problem}[Erd\H{o}s \#68]\label{long68:res:problem}
Is $\Ser=\sum_{n\ge2}(n!-1)^{-1}$ irrational?
\end{problem}

\noindent Erd\H{o}s states the question on p.~102 of his 1988 survey and, in
the same passage, records the expectation that $\sum_n1/(n!+t)$ is
irrational, indeed transcendental, for every integer $t$
\cite[p.~102]{erdos1988}.  That is conjectural context.  Numbering and
current status follow \href{https://www.erdosproblems.com/68}{Bloom's
Erd\H{o}s problem catalogue} \cite{bloom}; the problem is open.

Put
\[
 d_n=n!-1,\qquad H_M=\sum_{n=2}^{M}\frac1{d_n},\qquad
 \Lc_M=\lcm(d_2,\ldots,d_M),\qquad
 A_M=\sum_{n=2}^{M}\frac{\Lc_M}{d_n},
\]
so that $H_M=A_M/\Lc_M$.  Throughout, $\den(x)$ is the positive denominator
of a rational number in lowest terms.  The first question is which prime
power present in $\Lc_M$ remains in $\den(H_M)$.

\label{long68:res:lead-prime-pole}\phantomsection
\begin{theorem}[maximal prime-power survival]\label{long68:res:prime-pole}
Let $M\ge2$, let $p$ be a prime dividing $\Lc_M$, and put $e=v_p(\Lc_M)$.
Let $J=\{n:2\le n\le M,\ v_p(d_n)=e\}$ and write $d_n=p^eu_n$ for $n\in J$.
Then, with inverses in $\mathbb F_p$,
\begin{equation}\label{long68:eq:prime-pole-survival}
 v_p\bigl(\den(H_M)\bigr)=e
 \quad\Longleftrightarrow\quad
 \sum_{n\in J}u_n^{-1}\ne0\quad\hbox{in }\mathbb F_p.
\end{equation}
\end{theorem}

\begin{proof}
Write $\Lc_M=p^eW$ with $p\nmid W$.  If $v_p(d_n)<e$, then $p\mid\Lc_M/d_n$.
If $n\in J$, then $u_n(\Lc_M/d_n)=W$, so $\Lc_M/d_n\equiv Wu_n^{-1}\pmod p$.
Summing over $2\le n\le M$ gives
\begin{equation}\label{long68:eq:prime-pole-residue}
 A_M\equiv \frac{\Lc_M}{p^e}\sum_{n\in J}u_n^{-1}\pmod p.
\end{equation}
  The historical manuscript cites \path{Erdos68/PaperCompletePrimePole.lean}, line 118. This module is absent from the retained source pin; this citation does not establish kernel verification at that pin.

Finally $\den(H_M)=\Lc_M/\gcd(A_M,\Lc_M)$, whose $p$-valuation is $e$ exactly
when $p\nmid A_M$.  Since $W$ is invertible modulo $p$, that is exactly the
stated condition.
\end{proof}

The criterion decides the survival of the full exponent $e$.  A zero residue
gives $v_p(\den(H_M))<e$, and further congruences are needed to determine the
surviving exponent.  The proof uses no property of factorials beyond the
displayed denominators, so it applies to any finite sum of reciprocals.
It is the first valuation layer of the reciprocal-sum formula of Louwsma
and Martino \cite[Lemma~4.1, p.~10]{louwsma-martino}.  What is specific here is the behaviour of
the actual family $d_n=n!-1$, and that behaviour is not generic.

\paragraph{Two complete cancellations.}
Take $M=p-1$.  The indices below with $p\mid d_n$ are all of them, and every
one has valuation exactly one.
\begin{center}
\small
\begin{tabular}{rlll}
\toprule
\papertablehead
$p$ & indices $n$ & $d_n/p\pmod p$ & inverses modulo $p$\\
\midrule
$139$ & $69,\ 122,\ 137$ & $6,\ 49,\ 73$ & $116,\ 122,\ 40$\\
$2593$ & $349,\ 2243,\ 2591$ & $1508,\ 1566,\ 1678$ & $1367,\ 356,\ 870$\\
\bottomrule
\end{tabular}
\end{center}
The inverse sums are $278=2\cdot139$ and $2593$.  By
Theorem~\ref{long68:res:prime-pole}, $139\nmid\den(H_{138})$ and
$2593\nmid\den(H_{2592})$: each prime divides the common denominator of its
prefix and vanishes from the reduced one.  Both hit lists and both valuations
are checked by the recurrence
\[
 r_1=1,\qquad r_n\equiv nr_{n-1}\pmod{p^2},\qquad 0\le r_n<p^2,
\]
run through $n=p-1$.  A hit is an index with $r_n\equiv1\pmod p$, its lifted
cofactor is $(r_n-1)/p$ modulo $p$, and valuation two would mean $r_n=1$,
which occurs at neither prime.  That enumeration is a finite calculation
separate from the theorem.

A prime is \emph{prefix-private at $m$} when it divides $d_m$ and divides
none of $d_2,\ldots,d_{m-1}$.

\begin{proposition}[cofinal first prime occurrences]\label{long68:res:wilson-cofinality}
For every integer $B\ge0$ there are a prime $q$ and an integer $m>B$ with
$m<q$, $q\mid m!-1$ and $\gcd(q,k!-1)=1$ for every $k$ with $2\le k<m$.
\end{proposition}

\begin{proof}
Choose a prime $q\ge B!+5$.  Wilson's theorem gives $(q-2)!\equiv1\pmod q$,
so the set of indices $n\ge2$ with $q\mid n!-1$ is nonempty and contains no
element above $q-2$; let $m$ be its least element.  If $m\le B$, then
$q\le m!-1\le B!-1$, contradicting the choice of $q$.  Hence $m>B$, and
minimality gives the coprimality assertions.
\end{proof}

The construction gives no useful upper bound for $q$ in terms of its first
hit, and that comparison is what the scale inequalities of
\S\ref{long68:sec:residue} consume.  Wilson reflection limits what can be inferred
from a large prime factor alone: for odd $n<q$ with $q\mid n!-1$, the
identity $(q-1-n)!\,n!\equiv(-1)^{n+1}\pmod q$ gives a second hit at
$q-n-1$.  The reflection identity is classical; see Stewart
\cite[p.~462, (4)]{stewart2004}.

\statusboundary

\section{The growth of the common denominator}\label{long68:sec:lcm}

Theorem~\ref{long68:res:prime-pole} concerns the denominator after reduction.  The
common denominator before reduction admits an unconditional lower bound of
its own, by an elementary argument.

\begin{lemma}[product, least common multiple, pairwise gcd]
\label{long68:res:product-lcm}
For positive integers $x_1,\ldots,x_k$,
\[
 \prod_{i=1}^{k}x_i\ \Big|\ \lcm(x_1,\ldots,x_k)\prod_{i<j}\gcd(x_i,x_j).
\]
\end{lemma}

\begin{proof}
Fix a prime $r$ and relabel so that $a_1\le\cdots\le a_k$, where
$a_i=v_r(x_i)$.  The right-hand side has $r$-valuation
$a_k+\sum_{i<j}\min(a_i,a_j)=a_k+\sum_{i=1}^{k-1}(k-i)a_i$, and the left-hand
side has $\sum_{i=1}^{k}a_i$.  The difference is
$\sum_{i=1}^{k-1}(k-i-1)a_i$, which is nonnegative.
\end{proof}

\begin{lemma}[factorial-gap gcd]\label{long68:res:gap-gcd}
For $2\le i<j$, the integer $g=\gcd(i!-1,j!-1)$ divides $j!/i!-1$, and
$g\le j!/i!-1<j^{\,j-i}$.
\end{lemma}

\begin{proof}
Both $i!$ and $j!$ are congruent to $1$ modulo $g$.  The quotient
$j!/i!=(i+1)(i+2)\cdots j$ is an integer, and $j!=i!\cdot(j!/i!)$, so
$j!/i!\equiv1\pmod g$.  Since $j!/i!\ge i+1>1$, the positive integer
$j!/i!-1$ is a multiple of $g$, whence $g\le j!/i!-1$.  Finally
$j!/i!$ is a product of $j-i$ integers each at most $j$.
\end{proof}

\begin{lemma}[segment inequality]\label{long68:res:segment}
For $2\le k\le N-1$,
\begin{equation}\label{long68:eq:segment}
 \sum_{n=N-k+1}^{N}\log(n!-1)
 \ \le\ \log \Lc_N+\binom{k+1}{3}\log N.
\end{equation}
  The historical manuscript cites \path{Erdos68/PaperCompleteGcdSegment.lean}, line 28. This module is absent from the retained source pin; this citation does not establish kernel verification at that pin.

\end{lemma}

\begin{proof}
Apply Lemma~\ref{long68:res:product-lcm} to $x_n=d_n$ for $N-k+1\le n\le N$.  Their
least common multiple divides $\Lc_N$.  By Lemma~\ref{long68:res:gap-gcd} each
pairwise gcd is smaller than $N^{\,j-i}$, and over a block of $k$ consecutive
indices
\[
 \sum_{i<j}(j-i)=\sum_{d=1}^{k-1}d(k-d)=\binom{k+1}{3}.
\]
Taking logarithms of the resulting divisibility gives \eqref{long68:eq:segment}.
\end{proof}

\begin{theorem}[common-denominator growth]\label{long68:res:lcm-growth}
\[
 \liminf_{N\to\infty}\frac{\log \Lc_N}{N^{3/2}\log N}
 \ \ge\ \frac{2\sqrt2}{3}.
\]
  The historical manuscript cites \path{Erdos68/PaperCompleteLiminf.lean}, line 42. This module is absent from the retained source pin; this citation does not establish kernel verification at that pin.
\end{theorem}

\begin{proof}
Fix $\alpha>0$ and take $k=\lfloor\alpha\sqrt N\rfloor$, which satisfies
$2\le k\le N-1$ for all large $N$.  Put $u=N-k+1$.

For the left-hand side of \eqref{long68:eq:segment}, use $n!-1\ge n!/2$ and
$\log n!\ge n\log n-n$.  Since $n\mapsto n\log n-n$ increases,
\[
 \sum_{n=u}^{N}\log(n!-1)\ \ge\ k\,(u\log u-u)-k\log2 .
\]
Here $u\ge N-\alpha\sqrt N$ and $k=\alpha\sqrt N+O(1)$, so
$k\,u\log u=\alpha N^{3/2}\log N+O(N\log N)$, while $k\,u=O(N^{3/2})$ and
$k\log2=O(\sqrt N)$.  Hence the left-hand side is
$\alpha N^{3/2}\log N+o(N^{3/2}\log N)$.

For the right-hand side, $\binom{k+1}{3}\le(k+1)^3/6$, so
$\binom{k+1}{3}\log N=\tfrac{\alpha^3}{6}N^{3/2}\log N+o(N^{3/2}\log N)$.
Therefore
\[
 \liminf_{N\to\infty}\frac{\log \Lc_N}{N^{3/2}\log N}
 \ \ge\ \alpha-\frac{\alpha^3}{6}.
\]
The right-hand side is maximised at $\alpha=\sqrt2$, with value
$\sqrt2-\tfrac{2\sqrt2}{6}=\tfrac{2\sqrt2}{3}$.
\end{proof}

The proof is an ordinary proof and is not kernel-checked.  Its two
ingredients are Lemmas~\ref{long68:res:product-lcm} and~\ref{long68:res:gap-gcd}, and it
needs no factorial-congruence multiplicity theorem.  A weaker exponent
$4/3$ follows instead from Theorem~12 of Garaev, Luca and Shparlinski
\cite[arXiv v1, Thm.~12, p.~16]{garaev-luca-shparlinski}, which bounds by
$O(N^{2/3})$ the number of solutions of $n!\equiv a\pmod r$ with
$r$ prime, $a\not\equiv0\pmod r$, and $H+1\le n\le H+N$ where
$0\le H<H+N<r$; that deduction is recorded in the complete reasoning record and
is superseded here.  No antecedent for a lower bound on the least common
multiple of the numbers $n!-1$ was located, and the searches were for
factorial values shifted by one; least common multiples of consecutive
integers and of linear recurrences are a different family.  The finding is
that the statement is possibly known and was not located.

The constant $2\sqrt2/3$ is the best this argument gives.  The gain from a
block of $k$ indices is largest at the top of the range and the pairwise-gcd
loss grows with the spread, so consecutive top indices are simultaneously
optimal for both terms, and the bound $j!/i!<N^{\,j-i}$ is sharp to first
order there.  Convergence is slow: the omitted term $-ku$ is of order
$N^{3/2}$, so the finite ratio falls short of the limit by about
$\sqrt2/\log N$.

Theorem~\ref{long68:res:lcm-growth} is a statement about $\Lc_N$ and not about
$\den(H_N)$.  Theorem~\ref{long68:res:prime-pole} and the cancellations at $139$ and
$2593$ show why the substitution fails: a prime may occupy the common
denominator and be absent from the reduced one.  Bounding the reduced
denominator from below at every index is a stronger assertion, and it is
exactly what an irrationality proof through the clearing-scale route needs.

\section{Carries and the rationality boundary}\label{long68:sec:carry}

Define the strict integer successor and the predecessor gap by
\[
 Z_m=\lfloor m!H_m\rfloor+1,\qquad
 \Delta_m=Z_{m-1}-(m-1)!H_{m-1}\quad(m\ge3),
\]
so $0<\Delta_m\le1$, and define the integer carry $b_m$ by
\begin{equation}\label{long68:eq:carry-recurrence}
 Z_m=mZ_{m-1}+1-b_m .
\end{equation}
  The historical manuscript cites \path{Erdos68/PaperCompleteGcdSegment.lean}, line 73. This module is absent from the retained source pin; this citation does not establish kernel verification at that pin.

Call $b_m=1$ a \emph{unit carry}.  Put
$E_m=m!(\Ser-H_m)$ and $\varepsilon_m=1/(m!-1)$.  Since
$d_{n+1}>(n+1)d_n$, a geometric majorant gives the tail estimate
\begin{equation}\label{long68:eq:tail-bound}
 0<E_m<\frac{2\,m!}{(m+1)!-1}<\frac2m\le1\qquad(m\ge2).
\end{equation}

\label{long68:res:lead-carry-equivalence}\phantomsection
\label{long68:res:strict-successor-complete-characterization}\phantomsection
\begin{theorem}[strict-successor characterisation]\label{long68:res:carry-equivalence}
For $m\ge3$,
\begin{equation}\label{long68:eq:unit-window}
 b_m=1
 \iff m\mid Z_m
 \iff 1+\varepsilon_m<m\Delta_m\le2+\varepsilon_m .
\end{equation}
  The historical manuscript cites \path{Erdos68/PaperCompleteExisting.lean}, line 79. This module is absent from the retained source pin; this citation does not establish kernel verification at that pin.

Moreover
\begin{align}
 \Ser\in\Q
 &\iff b_m=1\ \hbox{for all sufficiently large }m,
 \label{long68:eq:carry-rationality}\\
 \Ser\notin\Q
 &\iff \forall B\ \exists m>B:\ m\nmid Z_m .
 \label{long68:eq:strict-misses}
\end{align}
If $\Ser=a/q$ with $a\in\Z$, $q\ge1$ and $b_m\ne1$, then $q\nmid(m-1)!$ and
$q\ge m$.
\end{theorem}

\begin{proof}
From $m!H_m=mZ_{m-1}-m\Delta_m+1+\varepsilon_m$ one gets
$b_m=\lceil m\Delta_m-1-\varepsilon_m\rceil$ with $-1\le b_m\le m-1$, which
gives both equivalences in \eqref{long68:eq:unit-window}, endpoints included.

Suppose $\Ser=a/q$ and $q\mid(m-1)!$.  For $j=m-1$ and $j=m$ the number
$j!\,\Ser$ is an integer, and $j!H_j=j!\,\Ser-E_j$ with $0<E_j<1$ by
\eqref{long68:eq:tail-bound}, so $Z_j=j!\,\Ser$.  Substituting into
\eqref{long68:eq:carry-recurrence} forces $b_m=1$.  Every fixed $q$ divides
$(m-1)!$ eventually, which gives one direction of
\eqref{long68:eq:carry-rationality} and, by contraposition, the final assertion;
$q<m$ would imply $q\mid(m-1)!$.  Conversely, an eventual unit-carry tail
makes $Z_m/m!$ eventually constant, and $0<Z_m/m!-H_m\le1/m!$ with
$H_m\to\Ser$ identifies that constant as $\Ser$, which is then rational.
Negating \eqref{long68:eq:carry-rationality} gives \eqref{long68:eq:strict-misses}.
\end{proof}

The divisibility conclusion retains prime-power multiplicities, so it
constrains valuations, prime support included.  It does not require a
prime divisor of $q$ larger than $m-1$.

\subsection{Why one proposed window argument is circular}
\label{long68:sec:adjacent-unit-no-go}

One route to cofinal non-unit carries tries to exclude two consecutive unit
carries by extracting a prime-power or quotient-gcd obstruction from their
cleared width-one window.  The exact two-step recurrence shows what that
window can deliver.  Let $U_k$ and $V_k$ be the numerator and denominator of
the reduced predecessor gap and let $G_k>0$ be the exact transition
normaliser.  For $m\ge3$, the two unit carries are equivalent to an integer
offset $\Omega_m$ with $0<\Omega_m\le D_m$, the window denominator telescopes
independently of the carry values,
\[
 D_m=V_{m+2}G_{m+1}G_m=V_m(m!-1)\bigl((m+1)!-1\bigr),
\]
and under the pair assumption the offset factors to match,
$\Omega_m=U_{m+2}G_{m+1}G_m$.  Cancelling the common positive normaliser
reduces the complete window to
\[
 0<U_{m+2}\le V_{m+2},
\]
which every reduced positive gap in $(0,1]$ already satisfies.  For that
specified window, therefore, an obstruction obtained only after imposing the
carry pair rewrites a bound that holds anyway.  The calculation does not
disprove adjacent unit carries and does not exclude an independent
restriction on the transition numerators or their valuations, which would
need its own argument constraining the predecessor state before the pair is
assumed.

\subsection{The companion coordinate}

Let $C=\sum_{n\ge2}\bigl(n!(n!-1)\bigr)^{-1}$.  The identity
$1/(n!-1)=1/n!+1/(n!(n!-1))$ and absolute convergence give
\begin{equation}\label{long68:eq:companion-decomposition}
 \Ser=C+e-2 .
\end{equation}
The canonical factorial digits of a real $x$ are
$a_m(x)=\lfloor m!x\rfloor-m\lfloor(m-1)!x\rfloor\in\{0,\ldots,m-1\}$, and
$x$ is rational exactly when they vanish from some index on.  That criterion
is due to Cantor \cite{cantor1869} and is recorded in Galambos
\cite[Ch.~1]{galambos1976}.

\label{long68:res:lead-companion-orbit}\phantomsection
\begin{proposition}[companion-orbit criterion]\label{long68:res:companion-orbit}
\[
 \Ser\in\Q
 \quad\Longleftrightarrow\quad
 \lfloor m!C\rfloor\equiv-2\pmod m
 \quad\hbox{for all sufficiently large }m.
\]
\end{proposition}

\begin{proof}
Put $J_m=\sum_{k=0}^{m}m!/k!$.  Then $J_m$ is an integer, every summand with
$k<m$ is divisible by $m$ and the summand at $k=m$ is $1$, so
$J_m\equiv1\pmod m$; and $0<m!\,e-J_m<1$ for $m\ge2$.

Suppose $\Ser$ is rational.  For all large $m$ both $(m-1)!\,\Ser$ and
$m!\,\Ser$ are integers, so $m\mid m!\,\Ser$.  Multiplying
\eqref{long68:eq:companion-decomposition} by $m!$ and using
$\lfloor -m!\,e\rfloor=-J_m-1$ gives
$\lfloor m!C\rfloor=m!\,\Ser+2\,m!-J_m-1\equiv-2\pmod m$.

Conversely, assume the congruence from some index on.  Because
$0\le a_m(C)<m$, for $m\ge3$ the congruence is equivalent to $a_m(C)=m-2$.
Choose $N$ beyond the exceptional indices.  The canonical expansion gives
\[
 C=\frac{\lfloor N!C\rfloor}{N!}+\sum_{m>N}\frac{m-2}{m!},
\]
and the telescope $\sum_{m>N}(m-1)/m!=1/N!$ turns this into
\[
 \Ser=C+e-2
 =\frac{\lfloor N!C\rfloor+1}{N!}+\sum_{m=2}^{N}\frac1{m!},
\]
which is rational.
\end{proof}

\subsection{Escape from a smaller interval}

Write $\theta_m=\{m!\,\Ser\}$.  The exact lower interval has width
$E_m/m\in(0,2/m^2)$, so the target shrinks quadratically.

\begin{proposition}[lower-interval criterion]\label{long68:res:lower-escape}
\begin{equation}\label{long68:eq:lower-escape}
 \Ser\notin\Q
 \quad\Longleftrightarrow\quad
 \forall B\ \exists m>B:\ E_m\le m\theta_{m-1}.
\end{equation}
For $m\ge3$ the finite condition
\begin{equation}\label{long68:eq:finite-escape}
 m\Delta_m\le1+\varepsilon_m
 \quad\hbox{or}\quad
 1+\varepsilon_m+\frac2m\le m\Delta_m
\end{equation}
implies the escape inequality in \eqref{long68:eq:lower-escape}.  Cofinally many
instances of \eqref{long68:eq:finite-escape} therefore imply $\Ser\notin\Q$.
\end{proposition}

\begin{proof}
If $\Ser$ is rational, the remainders $\theta_{m-1}$ vanish eventually while
$E_m>0$, so the escape inequality fails eventually.  Conversely, suppose
$m\theta_{m-1}<E_m$ for every sufficiently large $m$.  By
\eqref{long68:eq:tail-bound} this gives $m\theta_{m-1}<1$, so $a_m(\Ser)=0$ and
$\theta_m=m\theta_{m-1}$.  A positive remainder grows past one under that
recurrence, so the remainders vanish and $\Ser$ is rational.

For the finite implication, suppose $m\theta_{m-1}<E_m$.  The tail recurrence
is $mE_{m-1}=1+\varepsilon_m+E_m$, so $E_{m-1}>E_m/m>\theta_{m-1}\ge0$, and
with $0<E_{m-1}<1$ the strict-successor identity gives
$\Delta_m=E_{m-1}-\theta_{m-1}$.  Hence
\[
 1+\varepsilon_m<m\Delta_m=1+\varepsilon_m+E_m-m\theta_{m-1}
 \le1+\varepsilon_m+E_m<1+\varepsilon_m+\frac2m,
\]
which contradicts \eqref{long68:eq:finite-escape}.
\end{proof}

Escape at a single index does not give a non-unit carry there.  The exact
classification is
\begin{equation}\label{long68:eq:escape-classification}
 E_m\le m\theta_{m-1}
 \quad\Longleftrightarrow\quad
 b_m\ne1\ \hbox{ or }\ a_m(\Ser)=m-1
 \qquad(m\ge3),
\end{equation}
  The historical manuscript cites \path{Erdos68/PaperCompleteExisting.lean}, line 104. This module is absent from the retained source pin; this citation does not establish kernel verification at that pin.

so the maximal-digit branch is available.  It is realised: at $m=52$ exact
rational arithmetic gives $b_{52}=1$ together with
$52\,\Delta_{52}>1+\varepsilon_{52}+2/52$, the margin being about
$0.5689674908$.  Since $52$ is composite, this settles nothing about
prime indices, and the cofinal statement in
Proposition~\ref{long68:res:lower-escape} does not identify escape with a non-unit
carry.

\subsection{The whole shifted family}

Erd\H{o}s asked about $\Ser$ inside a family, and the same coordinate covers
the family.  For an integer $t\ge-1$ put
$\Ser_t=\sum_{n\ge2}1/(n!+t)$ and $C_t=\sum_{n\ge2}1/\bigl(n!(n!+t)\bigr)$,
so that $\Ser_t=-tC_t+(e-2)$.

\begin{theorem}[uniform family boundary]\label{long68:res:shift-family}
For every integer $t\ge-1$, the series $\Ser_t$ is rational exactly when
\[
 \bigl\lceil t\,m!\,C_t\bigr\rceil\equiv2\pmod m
\]
for all sufficiently large $m$, and irrational exactly when that residue is
missed cofinally.  The member $t=-1$ is $\Ser$, and the member $t=0$ is
$e-2$, whose orbit is constantly $0$ and therefore misses the residue class
at every $m\ge3$, so the boundary returns the irrationality of $e$.
\end{theorem}

The $t=0$ specialisation recovers a classical fact and is included because it
is the one member of the family the criterion decides, which fixes the sense
in which the criterion is exact.  For every other $t$ the
producer of cofinal misses is missing, and the family is open.

\section{From denominator structure to a scale criterion}\label{long68:sec:residue}

For an integer $p\ge3$ put $I_p=\{2,\ldots,2p-1\}$ and $F_p=(p-1)!$, and
define
\begin{align*}
 D_p&=\lcm_{\substack{i,j\in I_p\\ i<j}}\gcd(d_i,d_j),
 &C_p&=\lcm(F_p,D_p),\\
 \Lc^{\mathrm{blk}}_p&=\lcm(F_p,d_2,\ldots,d_{2p-1}),
 &R_p&=\Lc^{\mathrm{blk}}_p/C_p .
\end{align*}
The normalised core is $\widetilde C_p=C_p/F_p$, so that
\begin{equation}\label{long68:eq:core-normalisation}
 \Lc^{\mathrm{blk}}_p=C_pR_p=F_p\widetilde C_pR_p .
\end{equation}
  The historical manuscript cites \path{Erdos68/PaperCompleteExisting.lean}, line 125. This module is absent from the retained source pin; this citation does not establish kernel verification at that pin.

The factorial base $F_p$ belongs inside every one of these quantities, and
omitting it changes the scale comparisons below by a factorial factor.  Put
\[
 T_p=\sum_{i\in I_p}\frac{\Lc^{\mathrm{blk}}_p}{d_i},\qquad
 \rho_p=(-T_p)\bmod R_p,\qquad
 K_p=2p^2(2p-1)! ,
\]
with $\rho_p$ the least nonnegative representative.  All of these depend only
on a finite prefix.

The core records every prime-power layer occurring in at least two gaps,
together with the base.  A prime dividing $R_p$ therefore has a unique
maximal-valuation gap above the base valuation, so the argument of
Theorem~\ref{long68:res:prime-pole} applies to it and gives $\gcd(T_p,R_p)=1$.  This
is where the first theorem returns.  Consequently, when $R_p>1$, the ratio
$T_p/R_p=C_pH_{2p-1}$ is not an integer and
\begin{equation}\label{long68:eq:gap-normalisation}
 \frac{\rho_p}{R_p}
 =\bigl\lceil C_pH_{2p-1}\bigr\rceil-C_pH_{2p-1}\in(0,1),
\end{equation}
a gap to the next integer.

\begin{theorem}[complementary-residue criterion]\label{long68:res:global-residue}
Suppose that for every $B$ there is a prime $p>B$ with $R_p>1$ and
\begin{equation}\label{long68:eq:global-scale}
 (2p+1)\Lc^{\mathrm{blk}}_p<K_p\rho_p .
\end{equation}
Then $\Ser$ is irrational.
\end{theorem}

\begin{proof}
Suppose $\Ser=a/q$ with $q\ge1$, and take a block with $p>\max(q,3)$
satisfying the hypotheses.  Then $q\mid F_p$, so
$U=\Lc^{\mathrm{blk}}_p(\Ser-H_{2p-1})$ is a positive integer, and
$\Lc^{\mathrm{blk}}_p\Ser=R_p(C_p\Ser)$ is a multiple of $R_p$; hence
$U\equiv-T_p\pmod{R_p}$ and $U\ge\rho_p$.  For the other side,
$1/(n!-1)<2/n!$ for $n\ge3$, and a geometric majorant gives
\[
 0<\Ser-H_{2p-1}<\frac{2}{(2p)!}\cdot\frac{2p+1}{2p}=\frac{2p+1}{K_p}.
\]
Therefore $K_p\rho_p\le K_pU<(2p+1)\Lc^{\mathrm{blk}}_p$, contrary to
\eqref{long68:eq:global-scale}.
\end{proof}

Cancelling $R_p$ in \eqref{long68:eq:global-scale} through
\eqref{long68:eq:core-normalisation} displays the quantitative issue directly:
\begin{equation}\label{long68:eq:core-gap-scale}
 (2p+1)C_p<K_p\frac{\rho_p}{R_p}.
\end{equation}
The private modulus has gone.  A supply of large private factors therefore
supplies neither a small collision core nor a lower bound for the normalised
gap $\rho_p/R_p$, which is at most one.  That is the exact reason
Proposition~\ref{long68:res:wilson-cofinality} does not close the argument.

\label{long68:res:lead-moving-factor-split}\phantomsection
\label{long68:res:moving-factor-scale-split}\phantomsection
The moving-factor criterion is a more specialised sufficient condition.  For
a prefix-private prime $q$ at $m\ge4$ the tailored block parameter is
$p=\lfloor m/2\rfloor+1$, which need not itself be prime.  For the factors
$1$ and $q$ of the private modulus the factor-pair floor is
$\min(\rho_p,R_p/q)$, and its scale inequality splits exactly:
\begin{equation}\label{long68:eq:factor-split}
 (2p+1)\Lc^{\mathrm{blk}}_p<K_p\min(\rho_p,R_p/q)
 \quad\Longleftrightarrow\quad
 \begin{cases}
 (2p+1)\Lc^{\mathrm{blk}}_p<K_p\rho_p,\\
 (2p+1)C_pq<K_p .
 \end{cases}
\end{equation}
  The historical manuscript cites \path{Erdos68/PaperCompleteSupplementary.lean}, line 59. This module is absent from the retained source pin; this citation does not establish kernel verification at that pin.

Testing both entries of the minimum and using
\eqref{long68:eq:core-normalisation} gives the equivalence, so the reduction has no
opaque floor premise left.  Cofinal instances of both inequalities imply
irrationality.  The single global condition of
Theorem~\ref{long68:res:global-residue} is separately sufficient, so the pair is one
sufficient route among several.

\label{long68:res:split-factor-normalized-collision}\phantomsection
\label{long68:bdry:fixed-owner-absorption}\phantomsection
One structural limitation is worth recording beside these criteria.  If the
owner index $n$ is fixed and $p>n!-1$, then $n!-1$ divides $(p-1)!$, so its
private quotient in the block at $p$ is exactly one.  Large private quotients
seen at small blocks, such as the factor $719$ owned at $n=6$, are therefore
finite-range phenomena, and the selected factors must escape with $p$.

\section{A limitation of finite-channel cancellation}\label{long68:sec:channels}

For a finitely supported integer vector $c=(c_i)$ on indices $i\ge2$, put
\[
 M(c)=\sum_i c_i\,i!,\qquad
 W_{d,i}=\frac{i!}{(d!)^{\lfloor i/d\rfloor}},\qquad
 \chn dc=\sum_i c_iW_{d,i}\quad(d\ge2).
\]
Each $W_{d,i}$ is an integer: writing $i=kd+r$ with $0\le r<d$, the quotient
$i!/\bigl((d!)^kr!\bigr)$ is a multinomial coefficient.  Since
$i!=(d!)^{\lfloor i/d\rfloor}W_{d,i}$ and $d!\equiv1\pmod{d!-1}$,
\begin{equation}\label{long68:eq:channel-congruence}
 \chn dc\equiv M(c)\pmod{d!-1}.
\end{equation}

\begin{theorem}[integral normal form]\label{long68:res:normalform}
For every finite integer support and every $d\ge2$ there is an integer $k$
with $\chn dc=M(c)+(d!-1)k$.
\end{theorem}

\begin{theorem}[quotient-band breakpoint]\label{long68:res:bandbreakpoint}
Let $d\ge2$ and $k\ge0$, and suppose every supported index $i$ satisfies
$kd\le i<(k+1)d$.  Then $M(c)=(d!)^k\chn dc$.  In particular, cancellation
in the first band $d\le i<2d$ forces $M(c)=0$; and if every supported index
is at least $d$ while $M(c)\ne0$ and $\chn dc=0$, then some supported index
is at least $2d$.
\end{theorem}

\begin{proof}
Within one band the quotient $\lfloor i/d\rfloor$ is the constant $k$, so
$i!=(d!)^kW_{d,i}$ for every supported index and the factor $(d!)^k$ comes
out of the sum.  The two consequences follow by taking $k=1$ and by
contraposition.
\end{proof}

\noindent The hard step is the constant quotient; no valuation estimate
enters.  A nonzero moment therefore cannot be hidden entirely below $2d$
while the $d$-channel cancels.  This is a finite-family obstruction: it
constructs no cancelling family and says nothing about simultaneous channels
or residual size.

Consequently a vanishing $d$-th channel forces $(d!-1)\mid M(c)$,
and annihilating every channel $2\le d\le D$ forces $\Lc_D\mid M(c)$, where
$\Lc_D=\lcm_{2\le d\le D}(d!-1)$ is the same quantity as in
\S\ref{long68:sec:lcm}.  Every zero-moment variation of the support changes a
normalised channel contribution by an integer only, so such variations cannot
manufacture an extra fractional cancellation coordinate.

\subsection{The exact low-channel lattice and residual class}

The integral divisor basis makes the preceding divisibility statement exact.
Let $U_n$ be the isolated $n$-channel unit, let $K_D$ be the canonical vector
with moment $L_D$ and channels $2,\ldots,D$ equal to zero, and write
$a_D=[e_1]K_D$ and $u_n=[e_1]U_n$.  Every vector with zero initial coordinate
and vanishing channels through $D$ has a unique expression
\begin{equation}\label{long68:eq:low-channel-basis}
 c=tK_D+\sum_{n>D}z_nU_n,
 \qquad M(c)=tL_D,
\end{equation}
with finitely many nonzero integers $z_n$.  Requiring support on $n\ge2$
adds precisely the scalar equation
\begin{equation}\label{long68:eq:low-channel-support}
 ta_D+\sum_{n>D}z_nu_n=0.
\end{equation}

\begin{theorem}[exact attainable moment ideal]\label{long68:res:moment-ideal}
Fix $D\ge2$ and a prime $p$ with $D/2<p\le D$.  Put
\[
 H=D(2p-1),\qquad
 G_D=\gcd\{|u_n|:D<n\le H\},\qquad
 \mu_D=L_D\frac{G_D}{\gcd(G_D,a_D)}.
\]
The moments of finite integer vectors supported on $n\ge2$ and cancelling all
channels $2,\ldots,D$ are exactly $\mu_D\Z$.  The integer $\mu_D$ is positive,
is independent of the eligible prime $p$, and is attained by a vector of
coefficient content one, hence by a primitive vector.
\end{theorem}

\begin{proof}
The finite-horizon theorem gives $G_D>0$ and identifies it with the gcd of the
whole tail $\{u_n:n>D\}$; moreover $H<2D^2$.  Thus the finite integer
combinations in \eqref{long68:eq:low-channel-support} form $G_D\Z$, so that the
equation is soluble exactly when $G_D\mid ta_D$.  Dividing by
$\gcd(G_D,a_D)$ shows that $t$ is a multiple of
$G_D/\gcd(G_D,a_D)$, and \eqref{long68:eq:low-channel-basis} gives the displayed
moment ideal.  Bezout coefficients attain its positive generator.  If an
attaining vector had a nontrivial common coefficient divisor, division by that
divisor would produce a smaller positive attainable moment.  Hence its content
is one.  The ideal itself is intrinsic, so two eligible primes give the same
positive generator.
\end{proof}

For
\[
 \mathcal R(c)=\sum_{d\ge2}\frac{\chn dc}{d!-1},\qquad
 H_D=\sum_{2\le d\le D}\frac1{d!-1},
\]
the same coordinates retain the real information that the moment ideal alone
does not see.

\begin{theorem}[full residual transparency]\label{long68:res:residual-transparency}
For the vector in \eqref{long68:eq:low-channel-basis},
\[
 \mathcal R\!\left(tK_D+\sum_{n>D}z_nU_n\right)
 =tL_D(\Ser-H_D)+\sum_{n>D}z_n.
\]
The residual series converges for every finite vector supported away from
index zero.  A zero-moment vector has integral residual, and two such vectors
with the same factorial moment have residuals differing by an integer.
\end{theorem}

\begin{proof}
For $d>D$, the $K_D$ channel is $L_D$, while each $U_n$ contributes one to its
isolated residual coordinate.  Termwise decomposition therefore gives the
factorial-gap tail $L_D(\Ser-H_D)$ plus the finite coordinate mass
$\sum z_n$.  The factorial-gap tail is summable and the coordinate correction
has finite support, which proves convergence.  In the full divisor-basis
coordinates the coefficient of the $e_1$ column equals the factorial moment;
when that moment is zero only the integral coordinate mass remains.  Subtracting
two equal-moment identities proves the last assertion.
\end{proof}

The exact Lean endpoint packages are
\href{https://github.com/wcook04/plectis-erdos/blob/f20f727ae44d6dc3b24036bf63118b8fe591912e/ErdosProblems/Erdos68/PaperCompleteMomentIdeal.lean#L242}{primitive attainment of the moment-ideal generator} and
\href{https://github.com/wcook04/plectis-erdos/blob/f20f727ae44d6dc3b24036bf63118b8fe591912e/ErdosProblems/Erdos68/PaperCompleteResidualIdentity.lean#L195}{the equal-moment residual class}.
They determine the finite low-channel lattice and its residual class modulo
$\Z$; they do not prove that the remaining real residual is nonzero along a
cofinal family.

\label{long68:res:lead-channel-radius}\phantomsection
\begin{theorem}[square-subsequence radius bound]\label{long68:res:channel-radius}
Let $t,M,R\in\N$ satisfy
\[
 t\ge2^{32},\qquad M>0,\qquad \Lc_{2t^2}\mid M,\qquad M<(R+1)!-1 .
\]
Then $3t^3<2(R+1)$.  Consequently no family satisfying these hypotheses for
all sufficiently large $t$ has $(R(t)+1)/t^3\le3/2$ eventually, and none has
$R(t)=o(t^3)$.
  The historical manuscript cites \path{Erdos68/PaperCompleteExisting.lean}, line 236. This module is absent from the retained source pin; this citation does not establish kernel verification at that pin.
\end{theorem}

\begin{proof}
Put $D=2t^2$, $k=2t$ and $u=D-k+1=2t^2-2t+1$.  Lemma~\ref{long68:res:segment} at
$N=D$, with $n!-1\ge n!/2$ and $\log n!\ge n\log n-n$, and with
$\Lc_D\le M<(R+1)!$, gives
\begin{equation}\label{long68:eq:radius-log}
 2t\bigl(u\log u-u-\log2\bigr)
 <(R+1)\log(R+1)+\binom{2t+1}{3}\log(2t^2).
\end{equation}
  The historical manuscript cites \path{Erdos68/PaperCompleteSupportedBands.lean}, line 25. This module is absent from the retained source pin; this citation does not establish kernel verification at that pin.

Suppose $R+1\le\tfrac32t^3$.  Since $t\ge16$ one has
$\tfrac{15}{8}t^2\le u\le2t^2$ and $\log u\ge2\log t$, so the left side of
\eqref{long68:eq:radius-log} is at least
$\tfrac{15}{2}t^3\log t-4t^3-2t\log2$.  Using
$\binom{2t+1}{3}<\tfrac43t^3$ and $\log\tfrac32<\log2$, the right side is at
most
\[
 \tfrac32t^3\bigl(\log2+3\log t\bigr)+\tfrac43t^3\bigl(\log2+2\log t\bigr).
\]
Dividing by $t^3$ and rearranging, \eqref{long68:eq:radius-log} would require
\[
 \tfrac13\log t<4+\frac{2\log2}{t^2}+\tfrac{17}6\log2 .
\]
But $\log t\ge32\log2$ and $\tfrac23<\log2<1$ give
$\tfrac13\log t-\tfrac{17}6\log2\ge\tfrac{47}6\log2>5$, while
$4+2\log2/t^2<5$.  This contradiction proves the finite inequality, and the
two sequence conclusions follow.
\end{proof}

The hypotheses concern one specified cancellation architecture.  An
application to $\Ser$ must still produce the positive moment and the
factorial-size bound from the series, and must show that the channels through
$2t^2$ vanish for the resulting family.  The theorem restricts that
construction and supplies no irrationality argument.

Theorem~\ref{long68:res:lcm-growth} raises the asymptotic constant in the same
estimate.

\begin{corollary}[asymptotic radius constant]\label{long68:res:radius-constant}
Let $M(t),R(t)$ satisfy $M(t)>0$, $\Lc_{2t^2}\mid M(t)$ and
$M(t)<(R(t)+1)!-1$ for all sufficiently large $t$.  Then
\[
 \liminf_{t\to\infty}\frac{R(t)+1}{t^3}\ \ge\ \frac{16}{9}.
\]
  The historical manuscript cites \path{Erdos68/PaperCompleteLiminf.lean}, line 54. This module is absent from the retained source pin; this citation does not establish kernel verification at that pin.
\end{corollary}

\begin{proof}
Put $r=R(t)+1$, so $\Lc_{2t^2}\le M(t)<r!$ and $\log \Lc_{2t^2}<r\log r$.
Theorem~\ref{long68:res:lcm-growth} at $N=2t^2$ gives, for every $\varepsilon>0$ and
all large $t$,
\[
 \log \Lc_{2t^2}
 \ \ge\ \Bigl(\tfrac{2\sqrt2}{3}-\varepsilon\Bigr)\,2\sqrt2\,t^3
   \bigl(2\log t+\log2\bigr)
 \ \ge\ \Bigl(\tfrac{16}{3}-\varepsilon'\Bigr)t^3\log t .
\]
Suppose $r\le ct^3$ for arbitrarily large $t$, with $c<16/9$ fixed.  On that
subsequence $r\log r\le ct^3(3\log t+\log c)=\bigl(3c+o(1)\bigr)t^3\log t$,
and $3c<16/3$ contradicts the lower bound for large $t$.
\end{proof}

This corollary is asymptotic, and it does not give the strict finite
inequality $R+1>\tfrac{16}9t^3$ at every large $t$.  The finite logarithmic
constraint of the channel argument is compatible with equality
$9(R+1)=16t^3$ for $t\ge4$, so that constraint alone cannot force a strict
finite endpoint at $16/9$.  The two statements are separate, and
Corollary~\ref{long68:res:radius-constant} rests on the ordinary proof of
Theorem~\ref{long68:res:lcm-growth}; the kernel-checked finite theorem does not
carry it.

The congruence \eqref{long68:eq:channel-congruence} also explains the prime-channel
corrector.

\begin{theorem}[two-term prime channel corrector]\label{long68:res:translator}
Let $p\ge3$ be prime and let $c_{p-1}=p$, $c_p=-1$, with every other coefficient
zero.  Then $M(c)=0$, $\chn pc=p!-1$, and $\chn dc=0$ for every $d\ge2$ with
$d\ne p$.
\end{theorem}

\begin{proof}
Since $p\ge3$, both indices $p-1$ and $p$ lie in the prescribed coefficient
domain $i\ge2$.  The moment is $p\,(p-1)!-p!=0$.  For $2\le d<p$ the quotient identity
$\lfloor(p-1)/d\rfloor=\lfloor p/d\rfloor$ holds, since $d\nmid p$, so the
two channel weights carry the same power of $d!$ and the channel evaluates to
$p\,(p-1)!/(d!)^{\lfloor p/d\rfloor}-p!/(d!)^{\lfloor p/d\rfloor}=0$.  For
$d>p$ both indices lie below $d$, so both weights are the plain factorials
and the same cancellation occurs.  At $d=p$ the weights are $(p-1)!$ and
$p!/p!=1$, giving $p\,(p-1)!-1=p!-1$.
\end{proof}

For $p=2$, the same auxiliary identities would require $c_1=2$.  They do
not define an admissible corrector on the domain $i\ge2$, since every vector
supported on that domain has $c_1=0$.

\noindent At zero cost in the moment this supplies a unit in the
$p$-channel: adding an integer multiple of the corrector to any candidate
kernel shifts the $p$-channel by multiples of $p!-1$ and leaves the moment
and every other channel fixed.  For every channel rank and every prescribed
cutoff there is a factorial-grid kernel together with a remote prime-corrector
pair entirely beyond that cutoff, with all requested low channels zero,
nonzero moment, and residual in $[-1/2,1/2]$.  What is missing is strict
nonvanishing: nothing here rules out the rounded residual being exactly zero,
and no cofinal family with a strictly nonzero rounded residual has been
produced.

\section{Finite denominator exclusions}\label{long68:sec:finite}
\label{long68:res:lead-denominator-exclusions}\phantomsection

Two finite computations constrain a hypothetical denominator $q$ in
$\Ser=a/q$, and they constrain different features of it:
\begin{equation}\label{long68:eq:finite-bounds}
 q\nmid299999!,\qquad q\ge2^{39990},\qquad q>10^{12040}.
\end{equation}
The size bounds hold for every integer $a$ and positive natural $q$ with
$\Ser=a/q$; no reducedness assumption is needed. They are
\href{https://github.com/wcook04/plectis-erdos/blob/25ef6245d15a47548c6926369ae8f1a0f0a14a80/ErdosProblems/Erdos68/PaperCompleteFiniteSizeCertificate.lean#L58}{kernel-checked in Lean}.
The factorial-grid exclusion $q\nmid299999!$ is supported separately by the
external GMP computation. It is not a conclusion of the linked Lean theorem;
the full factorial-grid certificate remains an outstanding formal obligation.

Neither implies the other.  A denominator with no prime factor above
$299999$ can still fail to divide $299999!$ by carrying one prime to a high
power, and a denominator far below $2^{39990}$ can still fail that
divisibility; a magnitude bound above $299999!$ would imply it, and no such
bound is available here.

The first exclusion comes from an exact interval carry census.  An
independently implemented GMP integer computation certifies all $299998$
carry cells for $3\le m\le300000$ with a scale of $2^{5025679}$ and $96$
guard bits, using no floating-point arithmetic; at every step the outward
interval is proved to lie inside one half-open unit cell before the carry is
recorded.  Its unit carries occur exactly at
\[
 52,\ 591,\ 1030,\ 1407,\ 1438,\ 2164,\ 4258,\ 10991,\ 21236,
\]
so $b_{300000}\ne1$.  Theorem~\ref{long68:res:carry-equivalence} then gives
$q\nmid299999!$ and in particular $q\ge300000$.  The receipt is
\path{verification/erdos68-strict-successor.json}, schema
\path{erdos68_strict_successor_exact_interval_v1}, payload digest
\path{2974e8629f5959bdfa2814b63a9522cc9bf46e9e6ab99aaff28a8e980e1a1ee0},
recording the parameters above, the digests of the driver and backend used on
its own run, the event-trace digest and the final enclosure.  A later supplied
provenance record matches the receipt and both current driver hashes, so the
previously outstanding source-current regeneration check is discharged at
that exact scope; the census was not rerun as part of the present paper review.
The driver
\path{scripts/check_erdos68_strict_successor.py} and the backend
\path{scripts/check_erdos68_strict_successor_gmp.cpp} accompany it, and their
current hashes match the supplied provenance.  The same driver replayed independently at
$m\le4000$ reproduces the unit-carry prefix $52,591,1030,1407,1438,2164$ with
a non-unit carry at the endpoint, which gives $q\ge4000$ on its own.  The
census is a separate computation from the Lean development, and the
implication it feeds is the kernel-checked one.

The second exclusion has a short integer description.  Put $D=2^{80000}$ and
let $N$ be the first integer with $N!-1>D$.  Define
\[
 \ell=\sum_{n=2}^{N-1}\left\lfloor\frac D{n!-1}\right\rfloor,
 \qquad
 u=\ell+(N-2)+\left\lfloor\frac{2D}{N!-1}\right\rfloor+1 .
\]
Each rounded prefix term loses less than one, and $(n+1)!-1>(n+1)(n!-1)$
makes the tail from $N$ smaller than $2/(N!-1)$, so $\ell/D<\Ser<u/D$.  Here
$N=7054$ and $u-\ell=7053$.  Apply the continued-fraction algorithm to both
endpoints at once, retaining a quotient only while the integer parts agree
and both remainders are nonzero, and reversing the endpoint order at each
inversion.  There are $23449$ common quotients, and the convergent
denominators
\[
 Q_{-2}=1,\quad Q_{-1}=0,\quad Q_j=a_jQ_{j-1}+Q_{j-2}
\]
satisfy $Q_{23448}\ge2^{39990}$.  Every rational in the open enclosure shares
this initial segment, so its reduced denominator is at least $Q_{23448}$; the
argument is about a common prefix and does not require the last convergent to
lie inside the enclosure.  A newer exact width-$7056$ Farey certificate places
the enclosure between neighbours with denominator sum $B+E$, and checks
$B+E>10^{12040}$.  This separate denominator-sum comparison supplies the
strict decimal bound in \eqref{long68:eq:finite-bounds}; it does not follow
from $2^{39990}$.  The older width-$7053$ computation described here remains
a distinct historical certificate.  The
\href{https://github.com/wcook04/plectis-erdos/blob/cd3bd6fe245867a435e15d503ccedd699c5d02e2/scripts/check_erdos68_continued_fraction.py}{integer
replay} and its
\href{https://github.com/wcook04/plectis-erdos/blob/cd3bd6fe245867a435e15d503ccedd699c5d02e2/verification/erdos68-continued-fraction.json}{receipt}
record scale $80000$ bits, bracket width $7053$, $23449$ certified partial
quotients, power-of-two exponent $39990$ and strict decimal exponent $12038$.

Both certificates are finite.  A refinement of either enlarges the excluded
range and returns a statement of the same shape, and neither excludes an
eventual unit-carry tail.  The continued-fraction method is classical, and no
antecedent for these particular numerical outputs was located.

\section{The remaining arithmetic inputs}\label{long68:sec:open}

The exact target is \eqref{long68:eq:strict-misses}, equivalently
\eqref{long68:eq:lower-escape}.  The routes below are sufficient conditions for it,
and their quantitative hypotheses are unproved.  None of them is an
equivalent reformulation of the problem.

\paragraph{Weighted collisions and complementary residues.}
On the tailored moving-factor blocks, the local half of
\eqref{long68:eq:factor-split} is
\begin{equation}\label{long68:eq:weighted-target}
 \log\widetilde C_p
 <\log\left(\frac{2p^2}{2p+1}\cdot
       \frac{\prod_{j=p}^{2p-1}j}{q}\right),
\end{equation}
and the global half is \eqref{long68:eq:global-scale}.  Both must hold on the same
cofinal family.  For a prime $r$ absent from the base $F_p$, the contribution
to $v_r(\widetilde C_p)$ counts the levels $e\ge1$ at which at least two gaps
in $I_p$ are divisible by $r^e$; for a prime dividing $F_p$ the base
valuation must first be removed, so the layer count is
$v_r(\widetilde C_p)=\#\{e\ge1:h_{r,\,e+v_r(F_p)}(p)>1\}$ with
$h_{r,e}(p)=\#\{i\in I_p:r^e\mid i!-1\}$.  A spacing bound
$h_{r,e}(p)(e+1)\le2p+e-2$ is available.  An unweighted count of collisions
does not establish \eqref{long68:eq:weighted-target}, and reflected hits may not be
discarded, since Wilson reflection produces genuine collision primes.
Theorem~12 of Garaev, Luca and Shparlinski
\cite[arXiv v1, Thm.~12, p.~16]{garaev-luca-shparlinski} is a multiplicity
bound and supplies no lower bound for the complementary residue in
\eqref{long68:eq:global-scale}.

\paragraph{Repeated-support valuation amplification.}
Write $\Delta_n=u_n/v_n$ in lowest terms, let $B_n$ be the part of $n!-1$
supported on primes occurring in an earlier gap, and put
\[
 A_n=\prod_{\substack{r\mid B_n\\ v_r(v_n)<v_r(n!-1)}}r^{\,v_r(n!-1)} .
\]
What is established is $A_n\mid v_{n+1}$ and, when $A_n>1$,
$u_{n+1}\not\equiv0\pmod{A_n}$.  A sufficient quantitative goal, at cofinally
many genuinely nonterminal indices $n$, is the following with $m=n+1$ and
$d=A_n$:
\begin{equation}\label{long68:eq:amplification-target}
 \bigl((m+2)m!-2\bigr)v_m
 \le m^2(m!-1)\,(u_m\bmod d).
\end{equation}
It implies the upper alternative of \eqref{long68:eq:finite-escape}.  A nonzero
least representative can equal one while the modulus is enormous, so
divisibility and logarithmic mass do not by themselves bound that
representative.  A lower bound for $\log A_n$, and an infinitude of square
lifts, are separate possible inputs, and each needs a further implication to
reach \eqref{long68:eq:amplification-target}.

\paragraph{Lower-interval escape.}
Prove \eqref{long68:eq:finite-escape} at arbitrarily large indices; restricting to
primes is sufficient.  The upper alternative is exactly
$\bigl((m+2)m!-2\bigr)v_m\le m^2(m!-1)u_m$.  Replacing $u_m$ by its least
nonnegative residue modulo a specified divisor of $v_m$ gives a stronger
sufficient test.  Congruence recurrences alone do not count: synthetic models
satisfy the available congruences while remaining in the unit-carry branch,
and a zero canonical digit is a different event from membership in this
Archimedean interval.

\paragraph{Doubled-prime branch failure.}
For an odd prime $p$, the carry range and $Z_{2p}=2pZ_{2p-1}+1-b_{2p}$ give
\begin{equation}\label{long68:eq:doubled-prime}
 p^2\mid Z_{2p}
 \quad\Longleftrightarrow\quad
 \begin{cases}
 b_{2p}=1\ \hbox{and}\ p\mid Z_{2p-1},\ \hbox{or}\\
 b_{2p}=1+p\ \hbox{and}\ p\mid2Z_{2p-1}-1 .
 \end{cases}
\end{equation}
  The historical manuscript cites \path{Erdos68/PaperCompleteSupplementary.lean}, line 127. This module is absent from the retained source pin; this citation does not establish kernel verification at that pin.

Reduction modulo $p$ forces $b_{2p}\equiv1\pmod p$, and the two displayed
values are the only possibilities in $[-1,2p-1]$; division by $p$ gives the
predecessor conditions.  For rational $\Ser=a/q$ and $p>q$ the tail estimate
gives $Z_{2p}=(2p)!\,\Ser$, which is divisible by $p^2$.  Ruling out both
branches for infinitely many $p$ therefore proves irrationality.  The two
conditions must be coupled; controlling only the predecessor residue or only
the carry does not meet the hypotheses.

\paragraph{A remote Cramer residual.}
Put $s_n=((n+2)!)^2$ and $i_j=(t+j)s_n$ for $0\le j\le n+1$.  Form the
integer matrix $A$ with moment row $(i_j!)_j$ and channel rows
$(W_{d,i_j})_j$ for $2\le d\le n+2$, let $c$ be the cofactor vector of the
moment row, and let $N_d$ be the determinant obtained by replacing that row
with $(W_{d,i_j})_j$.  Cofactor expansion gives $M(c)=\det A$,
$\chn dc=N_d$, and $\chn dc=0$ for $2\le d\le n+2$.  By
\eqref{long68:eq:channel-congruence} each $(\chn dc-M(c))/(d!-1)$ is an integer and
vanishes for $d>\max_ji_j$, so
\[
 \mathcal R_{n,t}:=\sum_{d>n+2}\frac{N_d}{d!-1}
 =\det(A)\,\Ser+K_{n,t},\qquad K_{n,t}\in\Z .
\]
Seek a family with $\min_ji_j=t\,s_n\to\infty$ and
$\mathcal R_{n,t}\notin\Z$.  Such a family proves irrationality: for a fixed
denominator $q$, every entry of the moment row is divisible by $q$ once the
support is sufficiently remote, so $q\mid\det A$.  Unbounded pairs $(n,t)$ do
not supply that remoteness, because $t=0$ anchors the least support index at
zero.  The eventual identity $N_d=\det A$ controls the tail and gives no
control over the finite intermediate block, which changes sign; any
gcd-of-minors certificate must specify a positive gcd, its rank hypotheses,
and its connection to the displayed residual.

\paragraph{Criteria from the literature that do not apply.}
Duverney's fast-series criteria assume two-sided quadratic denominator growth
$cu_n^2\le u_{n+1}\le c'u_n^2$, whereas $u_{n+1}/u_n^2\to0$ for
$u_n=n!-1$.  His printed Corollary~3.2 displays the signed series
$\sum_n(u_{n+1}/u_n^2-1)$, without absolute values; its summands tend to
$-1$ here, so it does not converge to a finite value
\cite[pp.~275, 285--287]{duverney}.  The theorem of Barreto, Kang, Kim,
Kova\v{c} and Zhang proves irrationality of $\sum_n1/a_n$ under a
rapid-growth hypothesis that fails for $a_n=n!-1$, since
$\log(n!-1)=O(n\log n)$ makes $(n!-1)^{1/\psi^n}\to1$ for every fixed
$\psi>1$ \cite[Thms.~2--3 and Rem.~4]{barreto-et-al}.  Failure of a
sufficient criterion is not a theorem about $\Ser$, and the useful content is
the target these papers identify: a low-height clearing subsequence, or
enough exact cancellation in the least common multiple.
Theorem~\ref{long68:res:lcm-growth} bounds that cancellation from one side and
Theorem~\ref{long68:res:prime-pole} describes it from the other.

\medskip
\noindent Erd\H{o}s~\#68 is open.  Every rational representation with
positive denominator has $q\ge300000$.

\section*{Sources and evidence}

The formal counterparts below are checked by the Lean~4 kernel against
Mathlib at source snapshot \texttt{\commitshort}, with no \texttt{sorry},
added axioms, or unchecked evaluation.  Theorem~\ref{long68:res:lcm-growth} and
Corollary~\ref{long68:res:radius-constant} are ordinary proofs given in full above
and are not kernel-checked.  The two prefix cancellations, the index-$52$
example and the continued-fraction enclosure are finite integer calculations
with the procedures displayed above.  The carry census through $300000$ is a
separate exact-interval computation with the receipt named in
\S\ref{long68:sec:finite}.

\begin{center}\small
\begin{tabular}{@{}p{.46\linewidth}p{.48\linewidth}@{}}
\toprule
\papertablehead
Statement & Formal counterpart\\
\midrule
Theorem~\ref{long68:res:prime-pole} &
\lword{Erdos68/PrimePoleCriterion.lean}{223}{factorialGapPrefix_den_full_prime_power_iff}{maximal-power survival},
with the residue formula at
\lword{Erdos68/PrimePoleCriterion.lean}{129}{factorialGapPrefixLCMNumerator_mod_prime}{line 129}
and the denominator form at
\lword{Erdos68/PrimePoleDenominator.lean}{41}{natRatio_den_factorization_eq_iff}{line 41}\\
Proposition~\ref{long68:res:wilson-cofinality} &
\lword{Erdos68/PrimeZeroBranch.lean}{3097}{cofinal_prefixPrivate_factorialGap_hits}{cofinal private hits};
reflection at
\lword{Erdos68/PrimeZeroBranch.lean}{3139}{prime_dvd_reflected_factorialGap_of_odd}{line 3139}\\
Theorem~\ref{long68:res:carry-equivalence} &
\lword{Erdos68/FactorialGapPlateauCore.lean}{986}{irrational_factorialGapSeries_iff_cofinal_strictFacTopRat_misses}{strict-successor misses};
denominator exclusions at
\lword{Erdos68/FactorialGapPlateauCore.lean}{812}{rational_denominator_ge_of_nonunit_carry}{line 812}
and
\lword{Erdos68/FactorialGapPlateauCore.lean}{836}{rational_denominator_not_dvd_pred_factorial_of_nonunit_carry}{line 836}\\
Proposition~\ref{long68:res:companion-orbit} &
\lword{Erdos68/CompanionOrbitRationality.lean}{438}{not_irrational_factorialGapSeries_iff_eventually_companion_floor_neg_two}{companion orbit}\\
Proposition~\ref{long68:res:lower-escape} &
\lword{Erdos68/ShrinkingTargetNormalForm.lean}{82}{irrational_factorialGapSeries_iff_cofinal_lower_endpoint_escape}{lower-interval normal form};
finite window at
\lword{Erdos68/ShrinkingTargetNormalForm.lean}{122}{lower_endpoint_escape_of_predecessorGap_outside_window}{line 122},
cofinal consumer at
\lword{Erdos68/ShrinkingTargetNormalForm.lean}{145}{irrational_factorialGapSeries_of_cofinal_predecessorGap_outside_window}{line 145},
classification \eqref{long68:eq:escape-classification} at
\lword{Erdos68/ShrinkingTargetNormalForm.lean}{39}{factorialGap_lower_endpoint_escape_iff_nonunit_or_maximal_digit}{line 39},
radius at
\lword{Erdos68/ShrinkingTargetNormalForm.lean}{158}{lower_endpoint_target_radius_lt_two_div_sq}{line 158}\\
Theorem~\ref{long68:res:shift-family} &
\lword{Erdos68/FactorialShiftFamilyOrbit.lean}{139}{not_irrational_shiftGapSeries_iff_eventually_ceil_residue_two}{family boundary};
escape form at
\lword{Erdos68/FactorialShiftFamilyOrbit.lean}{156}{irrational_shiftGapSeries_iff_cofinal_ceil_residue_misses}{line 156},
member $t=-1$ at
\lword{Erdos68/FactorialShiftFamilyOrbit.lean}{175}{shiftGapSeries_neg_one_eq_factorialGapSeries}{line 175},
and the irrationality of $e$ at
\lword{Erdos68/FactorialShiftFamilyOrbit.lean}{211}{irrational_exp_one_of_family_boundary}{line 211}\\
Theorem~\ref{long68:res:global-residue} &
\lword{Erdos68/EndpointWeightedPrivateSupport.lean}{6164}{irrational_factorialGapSeries_of_cofinal_global_complementaryTail}{global complementary residue};
coprimality \eqref{long68:eq:gap-normalisation} at
\lword{Erdos68/EndpointWeightedPrivateSupport.lean}{4209}{factorialBlockTailNumerator_coprime_privateModulus}{line 4209}\\
Equation~\eqref{long68:eq:factor-split} &
\lword{Erdos68/EndpointWeightedPrivateSupport.lean}{4766}{factorialBlock_unitFactorPairFloor_scale_iff}{unit-factor scale split},
floor identity at
\lword{Erdos68/EndpointWeightedPrivateSupport.lean}{4701}{factorialBlock_unitFactorPairFloor_eq_min}{line 4701},
consumer at
\lword{Erdos68/PrimeZeroBranch.lean}{1574}{irrational_factorialGapSeries_of_cofinal_largePrefixPrivate_unitScaleSplit}{line 1574}\\
Fixed-owner absorption &
\lword{Erdos68/EndpointWeightedPrivateSupport.lean}{2677}{factorialBlockPrivateQuotient_eq_one_of_gap_lt}{absorbed owner}\\
Theorem~\ref{long68:res:normalform} and \eqref{long68:eq:channel-congruence} &
\lword{Erdos68/FactorialChannelCertificate.lean}{25}{factorial_pow_floor_dvd_factorial}{integral channel weight},
\lword{Erdos68/FactorialChannelCertificate.lean}{43}{channelWeight_mul_denominator}{exact cancellation}\\
Theorem~\ref{long68:res:channel-radius} &
\lword{Erdos68/ChannelIntegralCongruence.lean}{1222}{square_subsequence_radius_three_halves_lower}{finite radius bound};
sequence form at
\lword{Erdos68/ChannelIntegralCongruence.lean}{1241}{no_eventual_square_subsequence_three_halves_upper}{line 1241},
little-$o$ form at
\lword{Erdos68/ChannelIntegralCongruence.lean}{1043}{not_isLittleO_square_subsequence_radius}{line 1043},
method ceiling at
\lword{Erdos68/ChannelIntegralCongruence.lean}{764}{sharp_radius_satisfies_square_log_constraint}{line 764}\\
Prime channel corrector &
\lword{Erdos68/PrimeUnitTranslator.lean}{1657}{exists_remote_factorialGrid_primeTranslator_reduction}{remote reduction},
residual identity at
\lword{Erdos68/PrimeUnitTranslator.lean}{1559}{exists_cramerFactorialGridResidual_eq_det_mul_factorialGapSeries_add_int}{line 1559}\\
Equation~\eqref{long68:eq:doubled-prime} &
\lword{Erdos68/FactorialZeroPlateauSupplement.lean}{257}{sq_dvd_double_strictFacTop_factorialGapPrefix_iff}{doubled-prime criterion}\\
Canonical factorial digits &
\lword{Erdos68/CanonicalFactorialTermination.lean}{78}{exists_rat_eq_iff_eventually_zero_canonicalDigit}{termination equivalence}\\
Amplification modulus &
\lword{Erdos68/PrimeZeroBranch.lean}{5604}{factorialGapRepeatedBirthAmplificationModulus_dvd_succ_den}{divisibility},
nonvanishing at
\lword{Erdos68/PrimeZeroBranch.lean}{5619}{predecessorGapNumeratorNat_mod_repeatedBirthAmplificationModulus_ne_zero}{line 5619}\\
\bottomrule
\end{tabular}
\end{center}

Attribution.  Wilson's theorem and the Wilson reflection identity are
classical, the latter recorded by Stewart \cite[p.~462, (4)]{stewart2004}.
The factorial-digit termination criterion is Cantor's \cite{cantor1869}, in
the form recorded by Galambos \cite[Ch.~1]{galambos1976}; the additions here
are the telescope $C=\Ser-e+2$, the identification of the digit value $m-2$,
and the transport to the shifted family.  The multiplicity bound is Garaev,
Luca and Shparlinski's \cite{garaev-luca-shparlinski}, and the lcm deduction
from it is not theirs.  The survival criterion is a specialisation of
Louwsma and Martino's valuation formula \cite[Lemma~4.1]{louwsma-martino}.
For the Wilson cofinality packaging, the finite radius bound and
Theorem~\ref{long68:res:lcm-growth}, no antecedent was located, which is a negative search result, and no priority
is claimed.


\appendix
\section{Guide to the formal sources}\label{long68:app:sources}

The public \texttt{ErdosProblems.Erdos68} package carries the checked source
for this note.  The declarations below are the dependency chain behind the
statements above, together with the subsidiary results developed alongside
them; they are pinned to the formal-source commit named at the start of this
note.

\begin{itemize}[leftmargin=*]
\item \lref{Erdos68/CanonicalFactorialDigits.lean}{47}{facFloor_rat_eq_cleared}
\item \lref{Erdos68/CanonicalFactorialDigits.lean}{81}{canonicalDigit_eq_zero_of_rational}
\item \lref{Erdos68/CanonicalFactorialDigits.lean}{124}{canonicalDigit_eq_floor_mul_remainder}
\item \lref{Erdos68/CanonicalFactorialDigits.lean}{153}{canonicalDigit_lt_radix}
\item \lref{Erdos68/CanonicalFactorialDigits.lean}{162}{canonicalRemainder_recurrence}
\item \lref{Erdos68/CanonicalFactorialDigits.lean}{196}{factorial_expansion_partial}
\item \lref{Erdos68/CanonicalFactorialDigits.lean}{269}{zero_remainder_forces_zero_digit_tail}
\item \lref{Erdos68/CanonicalFactorialTermination.lean}{38}{canonicalRemainder_eq_zero_of_zero_digit_tail}
\item \lref{Erdos68/CanonicalFactorialTermination.lean}{54}{exists_rat_eq_of_canonicalRemainder_eq_zero}
\item \lref{Erdos68/FactorialGapPlateauCore.lean}{82}{factorialGrid_plateau}
\item \lref{Erdos68/FactorialGapPlateauCore.lean}{945}{irrational_factorialGapSeries_of_cofinal_nonunit_carries}
\item \lref{Erdos68/FactorialZeroPlateauCertificates.lean}{136}{sixty_seven_le_rational_denominator}
\item \lref{Erdos68/PrimeZeroBranch.lean}{6146}{factorialGap_zero_branch_iff_lower_endpoint_cylinder}
\item \lref{Erdos68/PrimeZeroBranch.lean}{3047}{exists_late_prefixPrivate_factorialGap_hit_of_primeProduct_lt}
\item \lref{Erdos68/PrimeZeroBranch.lean}{3231}{prime_dvd_factorialBlockNormalizedCollisionCore_of_odd_reflection}
\item \lref{Erdos68/PrimeZeroBranch.lean}{3252}{one_lt_factorialBlockPrimeHitCount_of_odd_reflection}
\item \lref{Erdos68/PrimeZeroBranch.lean}{1419}{factorialGapLargePrefixPrivatePowerModulus_dvd_tailoredBlockPrivateQuotient}
\item \lref{Erdos68/PrimeZeroBranch.lean}{1524}{factorialGapLargePrefixPrivatePrimes_unit_factor_pair_tailoredBlock}
\item \lref{Erdos68/AdjacentUnitCarryWindow.lean}{148}{consecutive_unit_carries_iff_positive_offset_le_den}
\item \lref{Erdos68/AdjacentUnitCarryWindow.lean}{215}{twoStep_den_mul_transitionNormalizers}
\item \lref{Erdos68/AdjacentUnitCarryWindow.lean}{243}{adjacentUnitCarryWindowDen_eq_twoStep_den}
\item \lref{Erdos68/AdjacentUnitCarryWindow.lean}{337}{adjacentUnitCarryWindowOffset_eq_twoStep_reducedNumerator}
\item \lref{Erdos68/ChannelBreakpointRigidity.lean}{71}{factorialMoment_eq_factorial_pow_mul_channelNumerator_band}
\item \lref{Erdos68/ChannelBreakpointRigidity.lean}{91}{factorialMoment_eq_zero_of_channelNumerator_eq_zero_band}
\item \lref{Erdos68/ChannelBreakpointRigidity.lean}{101}{factorialMoment_eq_factorial_mul_channelNumerator_firstBand}
\item \lref{Erdos68/ChannelBreakpointRigidity.lean}{130}{exists_index_ge_two_mul_of_factorialMoment_ne_zero_of_channel_eq_zero}
\item \lref{Erdos68/ChannelIntegralCongruence.lean}{1006}{square_subsequence_radius_cubic_lower}
\item \lref{Erdos68/ChannelIntegralCongruence.lean}{1023}{no_eventual_square_subsequence_cubic_upper}
\item \lref{Erdos68/FactorialChannelCertificate.lean}{65}{channelEvent_eq_zero_of_not_dvd}
\item \lref{Erdos68/FactorialChannelCertificate.lean}{156}{lambda34_residual_bounds}
\item \lref{Erdos68/FactorialAnalyticBoundary.lean}{52}{tendsto_factorialGap_succ_div_sq}
\item \lref{Erdos68/FactorialAnalyticBoundary.lean}{84}{not_summable_abs_factorialGap_succ_div_sq_sub_one}
\item \lref{Erdos68/FactorialAnalyticBoundary.lean}{131}{tendsto_factorialGap_rpow_one_div_two_pow}
\item \lref{Erdos68/EndpointWeightedPrivateSupport.lean}{556}{pairwiseCollisionCore_insert_gcd_lcm}
\item \lref{Erdos68/EndpointWeightedPrivateSupport.lean}{586}{collisionCore_insert_gcd_lcm}
\item \lref{Erdos68/EndpointWeightedPrivateSupport.lean}{698}{collisionCore_div_base_eq_pairwiseCollisionCore_div_gcd}
\item \lref{Erdos68/EndpointWeightedPrivateSupport.lean}{739}{collisionCore_div_base_factorization}
\item \lref{Erdos68/EndpointWeightedPrivateSupport.lean}{800}{collisionCore_div_base_mul_denominatorLcm_dvd_denominatorProd}
\item \lref{Erdos68/EndpointWeightedPrivateSupport.lean}{2707}{factorialBlockFixedPairPrivateQuotients_eq_one}
\item \lref{Erdos68/EndpointWeightedPrivateSupport.lean}{3327}{factorialBlockNormalizedCollisionCore_le_gapProd_div_gapLcm}
\item \lref{Erdos68/EndpointWeightedPrivateSupport.lean}{3394}{factorialBlockNormalizedCollisionCore_factorization}
\item \lref{Erdos68/EndpointWeightedPrivateSupport.lean}{3443}{factorialBlock_primePower_le_gapPow_of_two_hits}
\item \lref{Erdos68/EndpointWeightedPrivateSupport.lean}{3472}{factorialBlock_prime_hitPair_distance_gt_exponent}
\item \lref{Erdos68/EndpointWeightedPrivateSupport.lean}{3602}{factorialBlock_primePowerHitCount_mul_succ_le}
\item \lref{Erdos68/EndpointWeightedPrivateSupport.lean}{3790}{factorialBlock_normalizedCollision_factorization_add_base_lt_prime}
\item \lref{Erdos68/EndpointWeightedPrivateSupport.lean}{3848}{factorialBlock_normalizedCollision_factorization_add_base_lt_blockDiameter}
\item \lref{Erdos68/EndpointWeightedPrivateSupport.lean}{3982}{factorialBlock_prime_sq_gt_pred_of_dvd_normalizedCollisionCore}
\item \lref{Erdos68/EndpointWeightedPrivateSupport.lean}{4387}{factorialBlock_factorProjection_lcm_eq_privateModulus}
\item \lref{Erdos68/EndpointWeightedPrivateSupport.lean}{4854}{factorialBlock_complementaryFactorPairFloor_le_global}
\item \lref{Erdos68/EndpointWeightedPrivateSupport.lean}{5213}{factorialBlock_normalizedCollisionCore_factorization_eq_repeatedHitLayerCount}
\item \lref{Erdos68/EndpointWeightedPrivateSupport.lean}{5277}{factorialBlock_normalizedCollisionCore_factorization_eq_truncatedRepeatedHitLayerCount}
\item \lref{Erdos68/EndpointWeightedPrivateSupport.lean}{5389}{factorialBlock_primePower_dvd_normalizedCollisionCore_iff_one_lt_hitCount_of_endpoint_lt_base}
\end{itemize}

\clearpage
% ---- part extended_record ----
% Long-only authored mathematics for the Erdős #68 reasoning record.
% Everything here was in the short note before the September 2026 revision and
% is kept because the complete record owns the subsidiary arguments, the
% superseded deductions, and the failed routes with their exact boundaries.

\section{Extended record}\label{long68:sec:extended-record}

The short note above carries the advertised results with complete proofs.
This part keeps the surrounding mathematics: the digit kernel the carry and
orbit criteria run on, the exact interfaces of the collision-core reduction,
the finite certificates, the superseded deductions, and the routes that were
tried and closed.

\subsection{The canonical factorial digit kernel}\label{long68:sec:ext-digits}

Write $\theta_0=\{x\}$ and, for $m\ge1$,
\[
 a_m=\lfloor m\,\theta_{m-1}\rfloor,\qquad \theta_m=m\,\theta_{m-1}-a_m ,
\]
so that $\theta_m=\{m!\,x\}$ for $m\ge1$.  The kernel checks the floor
formula, the digit bounds $0\le a_m<m$, the recurrence
$\theta_{m+1}=(m+1)\theta_m-a_{m+1}$, the finite telescoping expansion
\[
 x=\lfloor x\rfloor+\sum_{m=2}^{N}\frac{a_m}{m!}+\frac{\theta_N}{N!},
\]
and the rule that a zero remainder at one index forces every later digit to
vanish.  The rational direction is also checked: if $q>0$ and $q\le n$, then
$\operatorname{facFloor}(a/q,n)=((n!/q):\Z)\,a$ and the canonical digit at
radix $n+1$ vanishes, so every rational input has an eventually zero
expansion.  The converse holds for every real input, and gives
\[
 x\in\Q\quad\Longleftrightarrow\quad a_m(x)=0\ \hbox{for all large }m .
\]
If all digits after index $N\ge1$ vanish, the recurrence gives
$\theta_{N+k}\ge(k+1)\theta_N$; since every remainder is below one,
$\theta_N=0$ and $x=\lfloor N!x\rfloor/N!$.

A zero canonical digit is a different event from a zero-branch hit.  The
returned data contain canonical zero digits at $m=5$ and $m=23$, while the
zero-branch list is empty through $m=100000$.

A second exact reformulation runs through a defect automaton.  For a rational
centre recurrence $F_m=mF_{m-1}+1+\varepsilon_m-C_m$, the kernel checks that
the integer ceiling defect code equals
$\lfloor m\delta_{m-1}-\varepsilon_m\rfloor$ and that
$\delta_m=m\delta_{m-1}-\varepsilon_m-q_m$, with the specialisation
$\varepsilon_m=1/(m!-1)$ written out.  What is checked is the algebra of the
automaton; proving that the finite-sum residual centre satisfies the premise
is a separate step and is not done.

\subsection{Criteria from the literature, in full}\label{long68:sec:ext-literature}

Koepf and Schmersau prove that eventual equality between the floors of $n$
times a partial sum and $n$ times its limit forces irrationality
\cite[Thm.~1.1, p.~117]{koepf-schmersau}; their rational-term version obtains
that equality from prefix integrality at a scale $p_n$ together with the
strict tail bound $a-s_n<1/(np_n)$
\cite[Thms.~2.2--2.3, pp.~119--120]{koepf-schmersau}.  For the natural
termwise clearing choice $p_n=\lcm\{k!-1:2\le k\le n\}$ the last two
denominators already obstruct it:
\[
 p_n\ \ge\ \frac{(n!-1)\bigl((n-1)!-1\bigr)}{n-1},
\]
because $\gcd(n!-1,(n-1)!-1)=\gcd((n-1)!-1,n-1)\le n-1$.  For $n\ge4$ the
first omitted summand $1/((n+1)!-1)$ then already exceeds $1/(np_n)$, so this
$p_n$ cannot satisfy their tail hypothesis.  Cancellation in the reduced
prefix denominator could give a smaller scale, and proving enough
cancellation is another form of the present denominator problem.

Duverney's Theorem~3.1 assumes two-sided quadratic denominator growth
$cu_n^2\le u_{n+1}\le c'u_n^2$, while $u_{n+1}/u_n^2\to0$ for $u_n=n!-1$
\cite[pp.~275, 285--286]{duverney}; his printed Corollary~3.2 displays the
signed series $\sum_n(u_{n+1}/u_n^2-1)$, without absolute values.  Its
summands tend to $-1$ here, so it does not converge to a finite value
\cite[Cor.~3.2, p.~287]{duverney}.  The ratio proof divides numerator
and denominator by $(n!)^2$; the resulting terms tend to zero while
$(1-1/n!)^2$ tends to one.

The theorem of Barreto, Kang, Kim, Kova\v{c} and Zhang proves irrationality
of $\sum_n1/a_n$ when $a_n^{1/2^n}\to\infty$, whereas $(n!-1)^{1/2^n}\to1$
\cite[Thms.~2--3, pp.~2--4]{barreto-et-al}; the limit is kernel-checked
through the bound $0\le\log(n!-1)\le n^2$.  Their proof identifies a useful
exact criterion: Mahler's elementary rationality floor is contradicted by
prefix-clearing integers $D_N$ whose cleared positive tails satisfy
$\liminf_ND_Nr_N=0$ \cite[Lem.~8 and Prop.~12, pp.~6, 9--12]{barreto-et-al}.
The ordinary product of the factorial-gap denominators is far too large for
that estimate, so a transfer needs a low-height clearing subsequence, or
enough exact cancellation in their least common multiple.

Dividing the strict-successor recurrence $Z_m=mZ_{m-1}+1-b_m$ by $m!$ and
telescoping gives the exact finite identity
\[
 \frac{Z_M}{M!}=\frac{Z_2}{2!}+\sum_{m=3}^{M}\frac{1-b_m}{m!},
\]
so the carry defects $1-b_m$ are genuine factorial-series coefficients.
Han\v{c}l and Tijdeman give exact rationality classifications for polynomial
coefficients and finite-difference criteria for broader ordinary factorial
series \cite[Thm.~3.1 and Cor.~3.1, pp.~390--391]{hancl-tijdeman}; their
denominator is the cumulative linear product $\prod_{n\le N}(an+b)$, and the
individual number $N!-1$ does not occur.  Applied to the display above, the Cantor--Oppenheim criterion
still needs $1-b_m\ne0$ infinitely often, which is the missing cofinal
assertion.

\subsection{A superseded deduction}\label{long68:sec:ext-superseded}

Before the elementary segment argument of the short note, the record derived
a weaker exponent from an external multiplicity theorem.  Put
$Q_N=\prod_{2\le n\le N}(n!-1)$.  For an odd prime $r$, let $m_r$ count the
indices $2\le n\le N$ with $r\mid n!-1$; such an index satisfies $n<r$.  The
factorial-congruence multiplicity estimate of Garaev, Luca and Shparlinski
\cite[arXiv v1, Thm.~12, p.~16]{garaev-luca-shparlinski}, applied on
$1\le n\le\min(N,r-1)$, gives $m_r\ll N^{2/3}$; the prime $2$ divides none of
these factors.  Writing $E_r=\max_{2\le n\le N}v_r(n!-1)$,
\[
 \log Q_N=\sum_r\sum_{n=2}^{N}v_r(n!-1)\log r
 \le\bigl(\max_r m_r\bigr)\sum_rE_r\log r
 \ll N^{2/3}\log \Lc_N ,
\]
and Stirling summation gives $\log Q_N\asymp N^2\log N$, whence
$\log \Lc_N\gg N^{4/3}\log N$.  Theorem~\ref{long68:res:lcm-growth} supersedes this:
its exponent is $3/2$, its constant is explicit, and it uses no external
theorem.  The deduction is retained because it is the shape the returned
analyses describe, and because it is the only place a multiplicity bound
enters the record.

The primitive lcm divisibility survives cofactor removal: factorial
valuations do not remove the channel obstruction once every common cofactor
divisor has been removed.  A corank-one cofactor and determinant construction
for such primitive kernels has not been formalised, and the divisibility
theorem does not establish it.

\subsection{Plateaux, first exit, and the finite denominator ladder}
\label{long68:sec:ext-plateau}

A second argument works on the rational grid in place of the channels.  Let
$H$ be a partial sum and $q$ a candidate denominator.  Writing $qH=k+r$ and
$q(\Ser-H)=u$, the next $q^{-1}$ grid point $(k+1)/q$ lies below $\Ser$
exactly when $1\le r+u$.  The factorial plateau theorem states that if
$H<G\le\Ser$, if $n!G$ is integral, and if $n!(\Ser-H)<1$, then the strict
successor of $n!H$ and the canonical floor of $n!\Ser$ are the same grid
integer.  Two rigidity statements follow: consecutive plateau floors, scaled
by the next radix, force the canonical digit to vanish; and any first-exit
offset $\delta\in[0,2)$ with carry $b=-\lfloor\delta\rfloor$ has
$b\in\{0,-1\}$, so the exit has exactly two alternatives.

The first-crossing argument continues to a denominator bound.  For a rational
grid level $G$ with first crossing $\tau$ by the literal partial sums, write
$G-H_{\tau-1}=a/v$ with $a,v>0$.  Then $v\ge\tau!-1$, and on the $-1$ exit
branch $v\ge\tau!(\tau!-1)$.  No coprimality hypothesis on $a$ and $v$ is
needed.

There is also a direct obstruction at prime indices.  With $\Delta_m$ as in
the short note, the kernel checks for $m\ge3$ the criterion
\[
 m\mid Z_m
 \quad\Longleftrightarrow\quad
 1+\frac1{m!-1}<m\Delta_m\le2+\frac1{m!-1},
\]
and if $\Ser=a/q$ with $q>0$, then $p\mid Z_p$ for every prime $p>q$, so one
exact missed prime $p$ gives $q\ge p$.  Exact kernel reduction gives
$11\nmid Z_{11}$, hence $q\ge11$; exact rational normalisation gives
$60\nmid Z_{60}$, $64\nmid Z_{64}$ and $67\nmid Z_{67}$, and the prime index
$67$ gives the checked bound
\[
 \Ser=\frac aq,\ q>0\quad\Longrightarrow\quad q\ge67 .
\]
The exact-interval census of the short note replaces $67$ by $300000$.

There is a finite peeling identity.  For $x\ne0,1$ and $K\ge0$,
\[
 \frac1{x-1}=\sum_{j=1}^{K}\frac1{x^j}+\frac1{x^K(x-1)} .
\]
When $x=k!$ and a chosen factorial scale is divisible by $(k!)^K$, the scaled
finite sum is integral and only the last term retains the factor $k!-1$ in
its denominator.  The identity isolates one residual fraction before exact
bounding, and it supplies no cofinal family of nonzero residuals.

One proposed strengthening is false and is recorded as false: the
divisibility $(m!-1)\mid\den(V_m)$ fails at the reported strict events
$m=52$ and $m=591$.  Only the Archimedean first-crossing lower bound
survives.

There is a second, more arithmetic mechanism at doubled prime indices, stated
as \eqref{long68:eq:doubled-prime} in the short note.  The formal theorem is not
restricted to hand-reduced indices; it is the general unique-slot prime-power
criterion specialised to the literal prefixes.

\subsection{The collision-core interfaces}\label{long68:sec:ext-collision}

Beyond the definitions used in the short note, the record keeps the exact
interfaces of the collision-core reduction.

The collision core has an exact incremental law.  For the positive
factorial-gap denominators, adjoining $d_a$ to a finite family $S$ gives
\[
 C(S\cup\{a\})=\lcm\!\Bigl(C(S),\ \gcd\bigl(d_a,\lcm_{j\in S}d_j\bigr)\Bigr),
\]
since finite-family gcd and lcm distributivity collapses the lcm of all
pairwise gcds against $d_a$ to a single gcd.  The same formula holds after
adjoining the distinguished base, so each step needs only the old denominator
lcm and the old core, with no pairwise rescan.

There is an exact product and lcm bound.  If $\widetilde C(S)$ is the core
after cancelling a positive distinguished base, $L(S)=\lcm_{j\in S}d_j$ and
$P(S)=\prod_{j\in S}d_j$, then $\widetilde C(S)L(S)\mid P(S)$, hence
$\widetilde C(S)\le P(S)/L(S)$.  For the factorial block this specialises to
\[
 \widetilde C_p\le\frac{\prod_{n\in I_p}(n!-1)}{\lcm_{n\in I_p}(n!-1)},
\]
an exact bridge from lower estimates for the factorial-gap lcm to upper
estimates for the normalised core.  It does not close the local scale bound.

The base cancellation is exact prime by prime.  With $F_p=(p-1)!$ and $C_p$
the unnormalised core,
\[
 \widetilde C_p=\frac{\lcm(F_p,C_p)}{F_p}=\frac{C_p}{\gcd(F_p,C_p)},\qquad
 v_r(\widetilde C_p)=v_r(C_p)-\min\{v_r(F_p),v_r(C_p)\} .
\]
So $r^e\mid\widetilde C_p$ exactly when the pairwise core carries
$r^{e+v_r(F_p)}$, which in the block forces two distinct gaps to be divisible
by that higher power.  The sharp surviving valuation cap
$v_r(\widetilde C_p)+v_r(F_p)<r$ then gives $p-1<r(r-1)<r^2$ for every
support prime, and $\widetilde C_p$ is coprime to $k!$ whenever
$k(k-1)\le p-1$.  This removes every factorial channel below the moving
square-root cutoff, and it does not bound the aggregate product of the
remaining large prime powers.

For collision estimates that already provide an upper-half hit, no exponent
is lost to normalisation: if $r$ divides a displayed gap at some $n\ge p$,
then $r\nmid F_p$ and $r^e\mid\widetilde C_p\iff r^e\mid C_p$ for every
$e>0$.  This has an exact incidence-count form,
\[
 r^e\mid\widetilde C_p
 \quad\Longleftrightarrow\quad
 1<\#\{i\in I_p:r^e\mid i!-1\},
\]
so a source estimate giving at most one $r^e$-hit deletes that exponent and
yields $v_r(\widetilde C_p)<e$.  The local aggregation is exact:
\[
 v_r(\widetilde C_p)=\#\bigl\{e\in[1,r-1]:1<\#\{i\in I_p:r^e\mid i!-1\}\bigr\},
\]
and for an endpoint prime with $2p-1<r$ the range truncates to
$e\in[1,2p-4]$.  The same equivalence holds under the exact condition
$r\nmid F_p$ in place of the endpoint inequality, which covers the entire
moving prime range at and above the block parameter; at $e=2$ it gives a
conditional squarefree conclusion $v_r(\widetilde C_p)\le1$.

The local load has a distance-sensitive witness.  If $r$ is prime, $e>0$ and
$r^e\mid\widetilde C_p$, there are $i<j$ in $I_p$ with
$r^{e+v_r(F_p)}\mid i!-1$, $r^{e+v_r(F_p)}\mid j!-1$ and
$r^{e+v_r(F_p)}\le j^{\,j-i}$; consequently $(2p-1)^d<r^{e+v_r(F_p)}$ forces
$d<j-i$.  The spacing hypothesis is discharged internally: any two $r^e$-hits
$i<j$ satisfy $e<j-i$, because $r\mid j!-1$ already forces $j<r$ and the
gap-power inequality converts that size relation into strict separation.  The
pairwise form is stronger than the selected-witness form: for arbitrary $r,e$
and any displayed hits $i<j$,
\[
 r^e\mid i!-1,\quad r^e\mid j!-1
 \quad\Longrightarrow\quad r^e\le j^{\,j-i},
\]
so every prime-power hit layer is an $e$-separated subset of the block and
\[
 (e+1)\,\#\{i\in I_p:r^e\mid i!-1\}\ \le\ 2p+e-2 .
\]
The unweighted packing step is therefore complete.  Globally,
\[
 r\mid\widetilde C_p
 \quad\Longrightarrow\quad
 v_r(\widetilde C_p)+v_r(F_p)<2p-3
\]
for every prime $r$ and $p\ge2$, so even primes already present in the
normalisation base pay for their base valuation inside the same
block-diameter budget.  What remains is to bound the layer counts uniformly
as $r$ and $p$ move, multiply the surviving contributions, and still close
the complementary-residue coordinate.

The squarefreeness premise is not proved.  An exhaustive modular scan through
$r\le2{,}000{,}000$ and $n\le240$ found four individual square hits and no
prime with two such hits.  All $498{,}501$ pairs $2\le a<b\le1000$ have
squarefree $\gcd(a!-1,b!-1)$, and the aggregate squarefree-collision scan
through $p=499$ stays below $0.374$ of the upper-descending-factorial
logarithmic scale.  These are finite exact computations.  They carry neither theorem authority
nor an asymptotic incidence bound.

A finite variant of the cofinal private supply compares the product of a
chosen set of primes, each at least $5$, with $\prod_{2\le k\le B}(k!-1)$; if
the prime product is larger, at least one chosen prime has no hit through
$B$, while Wilson still bounds its least hit by $q-2$.  Wilson reflection
limits what a linear-size divisor can be assumed to be: if $n$ is odd,
$n<q$, and $q\mid n!-1$, then $q\mid(q-n-1)!-1$, and when both indices lie in
one block with the reflected hit earlier, equivalently $q<2n+1$, the repeated
hit survives predecessor-factorial normalisation and its full-block incidence
count exceeds one.

Stewart states that for every $\varepsilon>0$ there are infinitely many odd
$n$ whose least prime factor $q$ of $n!-1$ satisfies
\[
 n<q<\left(\frac{\sqrt{145}-1}{8}+\varepsilon\right)n ,
\]
the printed text transferring estimate~(9) from $n!+1$ to $n!-1$
\cite[p.~464]{stewart2004}.  The source controls $q$ relative to the later
index $n$ and supplies no control relative to the private first-hit index,
so it is collision-core input.  The missing private-anchor and global product
estimates are elsewhere.

One selected prime $q$ already furnishes the coprime factor pair $(1,q)$: its
projection moduli are $R_p$ and $R_p/q$, whose least common multiple is
$R_p$, so no second selected prime is needed and there is no hidden
equality-or-disagreement branch in that specialisation.  The underlying
projection rigidity is general: if the endpoint numerator $Z$ is congruent to
a weighted numerator $T$ modulo $R$, then every divisor $Q$ of $R$ with
$Z\le B<Q$ satisfies $T\bmod Q=Z$, so two divisors above $B$ with unequal
projected residues exclude that bounded endpoint, and unequal projections
force $\min(Q_1,Q_2)\le T$.

\subsection{Finite channel and carry certificates}\label{long68:sec:ext-certificates}

The finite-support vector $\lambda=2e_3-e_4$ has, by kernel check,
$\chn2\lambda=0$, factorial moment $-12$, $\chn3\lambda=-2$,
$\chn4\lambda=11$, and $\chn d\lambda=-12$ for every $d\ge5$.  Under the
exact rational tail enclosure $1/119<\Theta_4<1/50$ its residual lies
strictly between $-93/575$ and $-309/13685$, so it is nonzero and subunit.

Exact integer regeneration verifies the canonical primitive kernels for every
$2\le D\le12$: channels $2$ through $D$ vanish, the factorial moment is
$\Lc_D$, and the coefficient content is one.  At $D=9$ the moment is
$\Lc_9=31540008254514077395$ and, after the stated prime-unit shift,
$1353/100000<R_9<1354/100000$.  At $D=3$ the vector $c=(-40,55,-10,1)$ on the
support $(3,4,5,6)$ annihilates channels $2$ and $3$, has moment $600$, and
satisfies
\[
 0.09925341997208298<\Lc_3(c)<0.09925341997208300 ,
\]
which excludes denominators dividing $600$.

There are two different computations at the same endpoint.  A returned
interval computation reports the stronger geometric statement that no
zero-branch event occurs at any $m\le100000$; its cited executable and source
digest were not supplied, so that classification remains external finite
evidence.  The strict-successor carry computation used in the short note is
local and independently regenerated, with its own source, backend, payload
and receipt digests.

\subsection{Limits of fixed-coordinate arguments}\label{long68:sec:ext-nogo}

\begin{itemize}
\item Residue vectors, their recurrences, and window widths admit synthetic
  all-hit blocks, so they cannot prove irrationality on their own.
\item Known pointwise prime congruences, prime-dilation congruences, parity,
  and the exact prime coefficient formula admit a synthetic rational
  countermodel.
\item Wilson quotients, harmonic sums, $p$-adic gamma identities, and
  factorial residues do not control the required Archimedean floor without an
  additional coupling theorem.  Every prime-window test factors into a sharp
  Archimedean strict-ceiling condition and a modular divisibility condition,
  and the missing ingredient is the coupling between them.
\item Fixed-denominator scalar canonical-product localisers and
  rank-saturated consecutive-jet Hermite--Pad\'e systems pay the full
  factorial-gap denominator.  For $E(z)=\prod_{n\ge2}(1-z/n!)$ the genus-zero
  product satisfies $-E'(1)/E(1)=\Ser$, and the natural scalar linear form
  carries the coefficient $Q_N=\prod_{2\le n\le N}(n!-1)$, for which $Q_N$
  times the tail diverges.  Exact first-order interpolation, scalar residue
  weighting, the natural Wronskian, and rank-saturated consecutive jets all
  reassemble the same prohibitive denominator.
\item Zero-moment variations cannot create an additional fractional
  cancellation coordinate, and factorial valuations cannot absorb the channel
  lcm obstruction.
\item A fixed pair of low-index private owners cannot make the projection
  argument cofinal, by the absorption statement of the short note.
\end{itemize}

\subsection{The generic shape behind the lower-interval form}
\label{long68:sec:ext-generic}

The cofinal statement in the short note has a generic ancestor.  For a real
$x$ and any positive sequence $c_m$ with $c_m<1/m$ eventually,
\[
 x\notin\Q
 \quad\Longleftrightarrow\quad
 \{(m-1)!\,x\}\ge c_m\ \hbox{for arbitrarily large }m .
\]
For rational $x$ the fractional parts vanish eventually; if all large
fractional parts lie below $c_m$, multiplying by $m$ stays below one, so the
next canonical digit is zero, and eventual zero digits force rationality.
The specific target $c_m=E_m/m$ is useful because it translates into finite
predecessor-gap conditions.  The smallness of the target alone establishes no
metric, measure or equidistribution theorem for this orbit, and none is
claimed.
% ---- part family_catalogue ----
\clearpage
\section{Complete result-family map}\label{sec:erdos-68-complete-family-map}

This section places every registered family for this problem in the shared 70-family reader order.  The five display bands control exposition only; the separate promotion state currently covers 6 families and is reported but does not hide or strengthen any family.  Mathematical statements and evidence modes come from the public claim registry.  Across all eight problems the public result-atom catalog contains 681 exact packet coordinates; this problem contributes 112.  The recovered source catalog contains 682 rows; 1 row(s) belong to families absent from the current claim registry and remain disclosed as detached source rows.  Catalog rows expose bounded statement excerpts plus full-source digests, not a claim that every complete packet statement is reproduced here.

\subsection{Factorial carry characterisation}
\textbf{Reader position.} 1 of 70; display band: front door.  Formal editorial disposition: promote.  These are separate classifications.

\textbf{Reader entry.} Exact equivalence between irrationality, cofinally many non-unit carries, and cofinal strict-successor divisibility failures.

Exact equivalence between irrationality, cofinally many non-unit carries, and cofinal strict-successor divisibility failures.

\textbf{Authority and reach.} Lean kernel; Comparator-selected; locally proved result; novelty unassessed.

\textbf{Exact boundary.} An exact reformulation does not supply the required cofinal failures.

\textbf{Result-atom population.} 33 of 681 public coordinates. The public atom catalog groups them under this family.

\textbf{Comparator assurance.} exact selected interface; 2 executable interface(s), 2 repository-registered selected result interface(s).

\textbf{Publication placement.} This family remains in the complete long record and is not a short-note headline.

\begin{itemize}
  \item \nolinkurl{ErdosProblems.Erdos68.irrational_factorialGapSeries_iff_cofinal_nonunit_carries}
  \item \nolinkurl{ErdosProblems.Erdos68.irrational_factorialGapSeries_iff_cofinal_strictFacTopRat_misses}
\end{itemize}

\subsection{Factorial channel and projection rigidity}
\textbf{Reader position.} 23 of 70; display band: mechanism.  Formal editorial disposition: split.  These are separate classifications.

\textbf{Reader entry.} Exact low-channel classification, the attainable moment ideal with a content-one primitive generator, full-residual transparency modulo integers, finite quotient-band factorisation and breakpoint, channel congruences, a two-term prime-channel corrector, and endpoint-weighted projection rigidity.

Exact low-channel classification, the attainable moment ideal with a content-one primitive generator, full-residual transparency modulo integers, finite quotient-band factorisation and breakpoint, channel congruences, a two-term prime-channel corrector, and endpoint-weighted projection rigidity.

\textbf{Authority and reach.} Lean kernel; locally proved result; novelty unassessed.

\textbf{Exact boundary.} These structural results determine the finite low-channel moment lattice and the residual class modulo integers, but they do not produce a cofinal nonintegral real residual or another obstruction settling Erdős 68.

\textbf{Result-atom population.} 63 of 681 public coordinates. The public atom catalog groups them under this family.

\textbf{Comparator assurance.} represented by selected interface; 0 executable interface(s), 0 repository-registered selected result interface(s).

\textbf{Publication placement.} This family is also admitted to the short note.

\begin{itemize}
  \item \nolinkurl{Erdos68.factorialMoment_eq_factorial_pow_mul_channelNumerator_band}
  \item \nolinkurl{Erdos68.exists_index_ge_two_mul_of_factorialMoment_ne_zero_of_channel_eq_zero}
  \item \nolinkurl{ErdosProblems.Erdos68.PaperComplete.attainable_moment_ideal}
  \item \nolinkurl{ErdosProblems.Erdos68.PaperComplete.exact_moment_ideal_with_primitive_attainment}
  \item \nolinkurl{ErdosProblems.Erdos68.PaperComplete.minimum_moment_content_one}
  \item \nolinkurl{ErdosProblems.Erdos68.PaperComplete.minimumMoment_independent_prime}
  \item \nolinkurl{ErdosProblems.Erdos68.PaperComplete.residual_transparency}
  \item \nolinkurl{ErdosProblems.Erdos68.PaperComplete.summable_fullResidual}
  \item \nolinkurl{ErdosProblems.Erdos68.PaperComplete.zero_moment_residual_integral}
  \item \nolinkurl{ErdosProblems.Erdos68.PaperComplete.equal_moment_residual_integer_difference}
  \item \nolinkurl{Erdos68.channelNumerator_mod_factorialMoment}
  \item \nolinkurl{Erdos68.primeTranslator_channelResidual_eq_one}
  \item \nolinkurl{ErdosProblems.Erdos68.projection_disagreement_excludes_bounded_endpoint}
\end{itemize}

\subsection{Factorial conditional producers}
\textbf{Reader position.} 43 of 70; display band: frontier.  Formal editorial disposition: hold.  These are separate classifications.

\textbf{Reader entry.} Cofinal zero-branch or large private-factor hypotheses imply the required irrationality conclusion.

Cofinal zero-branch or large private-factor hypotheses imply the required irrationality conclusion.

\textbf{Authority and reach.} Lean kernel; conditional reduction.

\textbf{Exact boundary.} The producer hypotheses are unproved.

\textbf{Result-atom population.} 1 of 681 public coordinates. The public atom catalog groups them under this family.

\textbf{Comparator assurance.} not selected for comparator; 0 executable interface(s), 0 repository-registered selected result interface(s).

\textbf{Publication placement.} This family is also admitted to the short note.

\subsection{Factorial lcm growth}
\textbf{Reader position.} 48 of 70; display band: frontier.  Formal editorial disposition: hold.  These are separate classifications.

\textbf{Reader entry.} A paper deduction gives a lower bound for the lcm of factorial-gap denominators using an external multiplicity theorem.

A paper deduction gives a lower bound for the lcm of factorial-gap denominators using an external multiplicity theorem.

\textbf{Authority and reach.} paper argument plus cited theorem; paper plus external theorem.

\textbf{Exact boundary.} Comparator cannot certify the cited input or authored deduction.

\textbf{Result-atom population.} 3 of 681 public coordinates. The public atom catalog groups them under this family.

\textbf{Comparator assurance.} not applicable to comparator; 0 executable interface(s), 0 repository-registered selected result interface(s).

\textbf{Publication placement.} This family is also admitted to the short note.

\subsection{Factorial finite certificates}
\textbf{Reader position.} 59 of 70; display band: technical support.  Formal editorial disposition: hold.  These are separate classifications.

\textbf{Reader entry.} Lean checks small exact misses; a hash-bound GMP scan reaches index 300000 and yields the corresponding finite denominator exclusion.

Lean checks small exact misses; a hash-bound GMP scan reaches index 300000 and yields the corresponding finite denominator exclusion.

\textbf{Authority and reach.} Lean finite theorem plus external exact certificate; finite computation.

\textbf{Exact boundary.} Finite computation does not change the cofinal quantifier.

\textbf{Result-atom population.} 12 of 681 public coordinates. The public atom catalog groups them under this family.

\textbf{Comparator assurance.} not applicable to comparator; 0 executable interface(s), 0 repository-registered selected result interface(s).

\textbf{Publication placement.} This family is also admitted to the short note.
% ---- part back ----
\begin{thebibliography}{9}
\small
\raggedright
\setlength{\itemsep}{0pt}
\bibitem{erdos1988} P.~Erd\H{o}s, \emph{On the irrationality of certain series:
  problems and results}, in A.~Baker (ed.), \emph{New Advances in Transcendence
  Theory}, Cambridge UP, 1988, pp.~102--109,
  doi:\href{https://doi.org/10.1017/CBO9780511897184.009}
  {10.1017/CBO9780511897184.009}.
\bibitem{bloom} T.~F. Bloom, \emph{Erd\H{o}s Problem \#68}.
  \url{https://www.erdosproblems.com/68}, accessed 28 July 2026.
\bibitem{louwsma-martino}
  J.~Louwsma and J.~Martino,
  \emph{Rational numbers with odd greedy expansion of fixed length},
  arXiv:\href{https://arxiv.org/abs/2309.07280v1}{2309.07280v1}, 2023,
  Lemma~4.1, p.~10.
\bibitem{cantor1869}
  G.~Cantor,
  \emph{\"Uber die einfachen Zahlensysteme},
  Z.\ Math.\ Phys.\ \textbf{14} (1869), 121--128.
\bibitem{galambos1976}
  J.~Galambos,
  \emph{Representations of Real Numbers by Infinite Series},
  Lecture Notes in Mathematics \textbf{502}, Springer, 1976, Chapter~1.
  doi:\href{https://doi.org/10.1007/BFb0081642}{10.1007/BFb0081642}.
\bibitem{garaev-luca-shparlinski}
  M.~Z. Garaev, F.~Luca, and I.~E. Shparlinski,
  \emph{Character sums and congruences with $n!$},
  Trans.\ Amer.\ Math.\ Soc.\ \textbf{356} (2004), no.~12, 5089--5102.
  \url{https://doi.org/10.1090/S0002-9947-04-03612-8};
  arXiv:\href{https://arxiv.org/abs/math/0403422}{math/0403422v1}.
\bibitem{stewart2004}
  C.~L. Stewart,
  \emph{On the greatest and least prime factors of $n!+1$, II},
  Publ.\ Math.\ Debrecen \textbf{65} (2004), no.~3--4, 461--480.
  \url{https://publi.math.unideb.hu/paper/989/download/10_5486_PMD_2004_3190.pdf}.
\bibitem{koepf-schmersau}
  W.~Koepf and D.~Schmersau,
  \emph{Irrationality of certain infinite series II},
  Analysis \textbf{31} (2011), 117--124.
  \url{https://doi.org/10.1524/anly.2011.1094}.
\bibitem{duverney}
  D.~Duverney,
  \emph{Irrationality of fast converging series of rational numbers},
  J.\ Math.\ Sci.\ Univ.\ Tokyo \textbf{8} (2001), 275--316.
  \url{https://www.ms.u-tokyo.ac.jp/journal/pdf/jms080206.pdf}.
\bibitem{hancl-tijdeman}
  J.~Han\v{c}l and R.~Tijdeman,
  \href{https://geodesic.mathdoc.fr/articles/10.4064/aa118-4-5/}
  {\emph{On the irrationality of factorial series}},
  Acta Arith.\ \textbf{118} (2005), no.~4, 383--401.
  doi:10.4064/aa118-4-5.
\bibitem{barreto-et-al}
  K.~Barreto, J.~Kang, S.-h.~Kim, V.~Kova\v{c}, and S.~Zhang,
  \emph{Irrationality of rapidly converging series: a problem of Erd\H{o}s
  and Graham},
  arXiv:\href{https://arxiv.org/abs/2601.21442v3}{2601.21442v3}, 2026,
  Theorems~2--3 and Remark~4.
\end{thebibliography}

\end{document}
